% ============================================================
%   CHAPITRE : Nombres complexes
%   Niveau   : Terminale S / Terminale C
%   Mazava Maths
% ============================================================
\chapter{L'ensemble des nombres complexes}

\begin{citationchapitre}
	« L'imaginaire est un beau et merveilleux refuge de l'esprit divin,\\
	presque un amphibie entre l'être et le non-être. »\\[0.4em]
	— Gottfried Wilhelm Leibniz
\end{citationchapitre}

\begin{objectifs}
	\begin{itemize}[label=\textbullet]
		\item Passer d'une forme à l'autre : algébrique, trigonométrique,
		exponentielle, polaire.
		\item Résoudre dans $\C$ les équations polynomiales de degré $2$, $3$
		et $4$, ainsi que les systèmes de deux équations à deux inconnues.
		\item \textbf{Interpréter géométriquement} module, argument, quotient
		d'affixes ; caractériser alignement, orthogonalité, cocyclicité,
		nature d'un triangle.
		\item \textbf{Déterminer et construire des ensembles de points} définis
		par une condition portant sur les affixes.
		\item Reconnaître l'écriture complexe des transformations usuelles.
	\end{itemize}
\end{objectifs}


% ============================================================
\section{Forme algébrique}
% ============================================================

\begin{definition}
	L'ensemble des \textbf{nombres complexes}, noté $\C$, est un ensemble contenant $\R$
	et un élément noté $\ii$ vérifiant $\ii^2 = -1$.

	Tout nombre complexe $z$ s'écrit de manière unique sous la forme
	$$
	z = a + b\ii, \quad a \in \R,\ b \in \R.
	$$
	Cette écriture s'appelle la \textbf{forme algébrique} de $z$.
	\begin{itemize}[label=\textbullet]
		\item $a = \mathrm{Re}(z)$ est la \textbf{partie réelle} de $z$.
		\item $b = \mathrm{Im}(z)$ est la \textbf{partie imaginaire} de $z$.
	\end{itemize}
\end{definition}

\remarque\\
$z$ est dit \textbf{réel} si $\mathrm{Im}(z) = 0$.\\
$z$ est dit \textbf{imaginaire pur} si $\mathrm{Re}(z) = 0$ et $z \neq 0$.

\begin{definition}[Égalité de deux complexes]
	$a + b\ii = c + d\ii \iff a = c \text{ et } b = d$.
\end{definition}

\begin{propriete}{Opérations en forme algébrique}{}

	Soient $z = a + b\ii$ et $z' = c + d\ii$ deux nombres complexes.
	\begin{itemize}[label=\textbullet]
		\item \textbf{Addition :} $z + z' = (a+c) + (b+d)\ii$
		\item \textbf{Multiplication :} $z \times z' = (ac - bd) + (ad + bc)\ii$
		\item \textbf{Conjugué :} $\bar{z} = a - b\ii$
	\end{itemize}
\end{propriete}

\begin{exemple}
	Soient $z = 3 + 2\ii$ et $z' = 1 - \ii$.
	\begin{itemize}[label=\textbullet]
		\item $z + z' = 4 + \ii$
		\item $z \times z' = (3 \times 1 - 2 \times (-1)) + (3 \times (-1) + 2 \times 1)\ii
		= 5 - \ii$
		\item $\bar{z} = 3 - 2\ii$
	\end{itemize}
\end{exemple}

\begin{propriete}{Propriétés du conjugué}{}
	Pour tous $z, z' \in \C$ :
	\begin{enumerate}
		\item $\overline{z + z'} = \bar{z} + \overline{z'}$
		\item $\overline{z \times z'} = \bar{z} \times \overline{z'}$
		\item $\overline{\left(\dfrac{z}{z'}\right)} = \dfrac{\bar{z}}{\overline{z'}}$ \quad ($z' \neq 0$)
		\item $z + \bar{z} = 2\,\mathrm{Re}(z) \in \R$ \quad et \quad
		$z - \bar{z} = 2\ii\,\mathrm{Im}(z)$
		\item $z \times \bar{z} = a^2 + b^2 \in \R^+$
		\item $\bar{\bar{z}} = z$
	\end{enumerate}
\end{propriete}

\begin{methode}
	\textbf{Calculer l'inverse ou le quotient d'un complexe :}
	Pour calculer $\dfrac{1}{z}$ ou $\dfrac{z}{z'}$, on multiplie numérateur et dénominateur
	par le \textbf{conjugué du dénominateur}.
\end{methode}

\begin{exemple}
	Calculer $\dfrac{3 + 2\ii}{1 - \ii}$.
	$$
	\frac{3 + 2\ii}{1 - \ii}
	= \frac{(3 + 2\ii)(1 + \ii)}{(1 - \ii)(1 + \ii)}
	= \frac{3 + 3\ii + 2\ii + 2\ii^2}{1 + 1}
	= \frac{3 - 2 + 5\ii}{2}
	= \frac{1}{2} + \frac{5}{2}\ii.
	$$
\end{exemple}


% ============================================================
\section{Module et représentation géométrique}
% ============================================================

\begin{definition}[Plan complexe]
	On munit le plan d'un repère orthonormé $(O, \vec{u}, \vec{v})$.
	À tout complexe $z = a + b\ii$, on associe le point $M(a\,;\,b)$
	appelé \textbf{image} de $z$, et on dit que $z=a+ib$ est l'\textbf{affixe} du vecteur $\vect{OM}$.
\end{definition}

\begin{definition}[Module]
	Le \textbf{module} de $z = a + b\ii$ est le réel positif
	$$
	\abs{z} = \sqrt{a^2 + b^2} = \sqrt{z \bar{z}}.
	$$
	Géométriquement, $\abs{z} = OM$ est la distance de $M$ à l'origine.
\end{definition}


% ============================================================
%  FIGURE : Module d'un nombre complexe
%  Utilisation de l'option "sloped" pour le label sur le vecteur
%  Mazava Maths — Chapitre Nombres complexes
% ============================================================

\begin{figure}[h]
	\centering
	\begin{tikzpicture}[>=Stealth, scale=1.1]

		% --- Grille -----------------------------------------------
		\draw[gray!30, thin] (0,0) grid (6,3);

		% --- Axes -------------------------------------------------
		\draw[->] (-0.3,0) -- (6.5,0)
		node[right, font=\small] {Axe des réels};
		\draw[->] (0,-0.3) -- (0,3.5)
		node[above, font=\small] {Axe des imaginaires};

		% --- Graduations axe réel ---------------------------------
		\foreach \x in {1,...,6}{
			\draw (\x,2pt) -- (\x,-2pt) node[below, font=\small] {\x};
		}
		\node[below left, font=\small] at (0,0) {$0$};

		% --- Graduations axe imaginaire ---------------------------
		\foreach \y in {1,2,3}{
			\draw (2pt,\y) -- (-2pt,\y) node[left, font=\small] {\y};
		}

		% --- Vecteurs de base u et v (en rouge) -------------------
		\draw[->, red, thick] (0,0) -- (1,0)
		node[below, font=\small, red] {$\vec{u}$};
		\draw[->, red, thick] (0,0) -- (0,1)
		node[left,  font=\small, red] {$\vec{v}$};

		% --- Tirets pointillés ------------------------------------
		\draw[blue!70, dashed, thin] (3.3,0) -- (3.3,2.2);   % vertical
		\draw[blue!70, dashed, thin] (0,2.2) -- (3.3,2.2);   % horizontal

		% --- Label b sur l'axe imaginaire -------------------------
		\node[left, font=\small, blue] at (0,2.2) {$b$};

		% --- Label a sur l'axe réel -------------------------------
		\node[below, font=\small, blue] at (3.3,0) {$a$};

		% --- Vecteur OM = module, label sloped --------------------
		\draw[->, red, thick] (0,0) -- (3.3,2.2)
		node[midway, above, sloped, font=\small, red]
		{$|z|=\sqrt{a^2+b^2}$};

		% --- Point M ----------------------------------------------
		\filldraw[blue] (3.3,2.2) circle (2pt)
		node[above right, font=\small, blue] {$M(a+ib)$};

	\end{tikzpicture}

	\smallskip
	\begin{minipage}{0.85\linewidth}
		\small
		Le \textbf{module} de $z = a+ib$ est la distance $OM=|z| = \sqrt{a^2 + b^2}$.
	\end{minipage}
	\caption{Représentation géométrique du module d'un nombre complexe.}
	\label{fig:module}
\end{figure}

M est un point d'affixe $z$ . \\
Alors le module de $z$ est égale à la distance $OM$.





\begin{exemple}
	$z = 3 + 4\ii$ : $\abs{z} = \sqrt{9 + 16} = \sqrt{25} = 5$.
\end{exemple}

\begin{propriete}{Propriétés du module}{}
	Pour tous $z, z' \in \C$ :
	\begin{enumerate}
		\item $\abs{z} \geqslant 0$ et $\abs{z} = 0 \iff z = 0$
		\item $\abs{\bar{z}} = \abs{z}$
		\item $\abs{z \times z'} = \abs{z} \times \abs{z'}$
		\item $\abs{\dfrac{z}{z'}} = \dfrac{\abs{z}}{\abs{z'}}$ \quad ($z' \neq 0$)
		\item $\abs{z^n} = \abs{z}^n$ \quad ($n \in \N$)
		\item $\abs{z + z'} \leqslant \abs{z} + \abs{z'}$ \quad (\textbf{inégalité triangulaire})
	\end{enumerate}
\end{propriete}

\begin{attention}
	En général, $\abs{z + z'} \neq \abs{z} + \abs{z'}$.
	Par exemple, $\abs{1 + (-1)} = 0 \neq \abs{1} + \abs{-1} = 2$.
\end{attention}

\interpretationgraphique
La distance entre les images de $z$ et $z'$ dans le plan complexe est $\abs{z' - z}$.


% ============================================================
\section{Forme trigonométrique et forme exponentielle}
% ============================================================

\begin{definition}[Argument d'un nombre complexe]
	Dans le plan complexe $(O, \vec{u}, \vec{v});$
	Soit $M$ un point d'affixe $z\neq 0$.

	On appelle \textbf{argument} de $z$, noté $\arg(z)$, toute mesure, en radians, de l'angle orienté $\left(\vec{u},\overrightarrow{OM}\right)$.
\end{definition}
% ============================================================
%  FIGURE : Module et argument d'un nombre complexe
%  Mazava Maths — Chapitre Nombres complexes
% ============================================================

\begin{figure}[h]
	\centering
	\begin{tikzpicture}[>=Stealth, scale=1.1]

		% --- Grille -----------------------------------------------
		\draw[gray!30, thin] (0,0) grid (5,3);

		% --- Axes -------------------------------------------------
		\draw[->] (-0.3,0) -- (5.8,0)
		node[right, font=\small] {Axe des réels};
		\draw[->] (0,-0.3) -- (0,3.5)
		node[above, font=\small] {Axe des imaginaires};

		% --- Graduations axe réel ---------------------------------
		\foreach \x in {1,...,5}{
			\draw (\x,2pt) -- (\x,-2pt) node[below, font=\small] {\x};
		}
		\node[below left, font=\small] at (0,0) {$0$};

		% --- Graduations axe imaginaire ---------------------------
		\foreach \y in {1,2,3}{
			\draw (2pt,\y) -- (-2pt,\y) node[left, font=\small] {\y};
		}

		% --- Vecteurs de base u et v (en rouge) -------------------
		\draw[->, red, thick] (0,0) -- (1,0)
		node[below, midway, font=\small, red] {$\vec{u}$};
		\draw[->, red, thick] (0,0) -- (0,1)
		node[left, midway, font=\small, red] {$\vec{v}$};

		% --- Tirets pointillés ------------------------------------
		\draw[gray!60, dashed, thin] (3,0) -- (3,2);
		\draw[gray!60, dashed, thin] (0,2) -- (3,2);

		% --- Arc de l'argument ------------------------------------
		% angle = arctan(2/3) ≈ 33.69°
		\draw[green!60!black, thick]
		(0.7,0) arc[start angle=0, end angle=33.69, radius=0.7];

		% --- Label arg(z) -----------------------------------------
		\node[green!60!black, font=\small] at (1.15,0.25)
		{$\arg(z)$};

		% --- Vecteur OM avec label |z| sloped ---------------------
		\draw[->, red, thick] (0,0) -- (3,2)
		node[midway, above, sloped, font=\small, red] {$|z|$};

		% --- Point M ----------------------------------------------
		\filldraw[blue] (3,2) circle (2pt)
		node[above right, font=\small, blue] {$M(z)$};

	\end{tikzpicture}

	\caption{Module et argument d'un nombre complexe $z$.}
	\label{fig:argument}
\end{figure}

	\smallskip
	\small
	Soit $z = a + ib$ de module $|z| = \sqrt{a^2+b^2}$. On pose $\theta = arg(z)$ tel que :
	$$
	\left\{
	\begin{array}{l}
		\cos\theta = \dfrac{a}{|z|}\\
		\sin\theta = \dfrac{b}{|z|}
	\end{array}
	\right.
	$$


	On écrit alors la \textbf{forme trigonométrique} :
	$$z = |z|\!\left(\cos\theta + i\sin\theta\right)$$.


\begin{definition}[Forme exponentielle]
	On note $\e^{\ii\theta} = \cos\theta + \ii\sin\theta$ (formule d'Euler).

	Tout complexe non nul s'écrit donc
	$$
	z = r\,\e^{\ii\theta}, \quad r = \abs{z} > 0,\quad \theta = \arg(z).
	$$
\end{definition}

\begin{exemple}
	\begin{itemize}[label=\textbullet]
		\item $z = 1 + \ii$ : $\abs{z} = \sqrt{2}$, $\arg(z) = \dfrac{\pi}{4}$,
		donc $z = \sqrt{2}\,\e^{\ii\pi/4}$.
		\item $z = -2$ : $\abs{z} = 2$, $\arg(z) = \pi$,
		donc $z = 2\,\e^{\ii\pi}$.
		\item $z = -\ii$ : $\abs{z} = 1$, $\arg(z) = -\dfrac{\pi}{2}$,
		donc $z = \e^{-\ii\pi/2}$.
	\end{itemize}
\end{exemple}

\begin{definition}[Forme polaire]
	La \textbf{forme polaire} de $z \neq 0$ est la donnée du couple
	$\bigl(r\,;\,\theta\bigr)$ où $r = \abs{z}$ et $\theta = \arg(z)$ ; on note
	$$
	z = [\,r\,;\,\theta\,].
	$$
	C'est une simple notation abrégée de $z = r(\cos\theta + \ii\sin\theta) = r\,\e^{\ii\theta}$.
\end{definition}

\begin{propriete}{Produit et quotient en forme polaire}{cplx_polaire}
	$$
	[\,r\,;\,\theta\,]\times[\,r'\,;\,\theta'\,] = [\,rr'\,;\,\theta+\theta'\,],
	\qquad
	\frac{[\,r\,;\,\theta\,]}{[\,r'\,;\,\theta'\,]} = \left[\,\frac{r}{r'}\,;\,\theta-\theta'\,\right].
	$$
	\emph{On multiplie les modules et on ajoute les arguments.}
\end{propriete}

\begin{propriete}{Récapitulatif des quatre formes}{cplx_formes}
	\begin{center}
	\renewcommand{\arraystretch}{1.4}
	\begin{tabular}{>{\raggedright\arraybackslash}p{2.6cm}
	                >{\raggedright\arraybackslash}p{5.2cm}}
	\hline
	\textbf{Forme} & \textbf{Écriture de $z$} \\
	\hline
	Algébrique & $z = a + b\ii$ \\
	Trigonométrique & $z = r(\cos\theta + \ii\sin\theta)$ \\
	Exponentielle & $z = r\,\e^{\ii\theta}$ \\
	Polaire & $z = [\,r\,;\,\theta\,]$ \\
	\hline
	\end{tabular}
	\end{center}
	avec $r = \sqrt{a^2+b^2}$, $\cos\theta = \dfrac{a}{r}$, $\sin\theta = \dfrac{b}{r}$,
	et réciproquement $a = r\cos\theta$, $b = r\sin\theta$.
\end{propriete}

\begin{propriete}{Propriétés des arguments}{}
	Pour tous $z, z' \in \C^*$ :
	\begin{enumerate}
		\item $\arg(z \times z') \equiv \arg(z) + \arg(z')\ [2\pi]$
		\item $\arg\!\left(\dfrac{z}{z'}\right) \equiv \arg(z) - \arg(z')\ [2\pi]$
		\item $\arg(z^n) \equiv n\,\arg(z)\ [2\pi]$ \quad ($n \in \Z$)
		\item $\arg(\bar{z}) \equiv -\arg(z)\ [2\pi]$
	\end{enumerate}
\end{propriete}

% ------------------------------------------------------------
\begin{resume}{Forme exponentielle}
	\label{res:cplx_expo}
	\textbf{Notation.} $\e^{\ii\theta} = \cos\theta + \ii\sin\theta$, donc tout
	$z \neq 0$ s'écrit $\boxed{z = r\,\e^{\ii\theta}}$ avec $r = \abs{z} > 0$ et
	$\theta = \arg(z)$.

	\smallskip
	\textbf{Calculs.} Les modules se \emph{multiplient}, les arguments s'\emph{ajoutent} :
	$$
	r\e^{\ii\theta}\times r'\e^{\ii\theta'} = rr'\,\e^{\ii(\theta+\theta')},
	\qquad
	\frac{r\e^{\ii\theta}}{r'\e^{\ii\theta'}} = \frac{r}{r'}\,\e^{\ii(\theta-\theta')},
	\qquad
	\bigl(r\e^{\ii\theta}\bigr)^{n} = r^{n}\,\e^{\ii n\theta},
	$$
	$$
	\overline{r\e^{\ii\theta}} = r\,\e^{-\ii\theta},
	\qquad
	\frac{1}{r\e^{\ii\theta}} = \frac{1}{r}\,\e^{-\ii\theta},
	\qquad
	-r\e^{\ii\theta} = r\,\e^{\ii(\theta+\pi)}.
	$$
	\textbf{Égalité.} $r\e^{\ii\theta} = r'\e^{\ii\theta'}
	\iff r = r'$ et $\theta \equiv \theta'\ [2\pi]$.

	\smallskip
	\begin{center}
		\begin{tikzpicture}[>=Stealth, scale=1.45]
			% axes
			\draw[->] (-1.35,0) -- (1.55,0) node[right, font=\scriptsize] {$\Re$};
			\draw[->] (0,-1.35) -- (0,1.5) node[above, font=\scriptsize] {$\Im$};
			% cercle unité
			\draw[black!65] (0,0) circle (1);
			% arc de l'argument
			\draw[green!60!black, thick]
			(0.32,0) arc[start angle=0, end angle=115, radius=0.32];
			\node[green!60!black, font=\scriptsize] at (0.60,0.33) {$\theta$};
			% vecteur générique
			\draw[->, red, thick] (0,0) -- (115:1);
			\node[red, font=\scriptsize, above left, inner sep=1pt] at (115:1) {$\e^{\ii\theta}$};
			% valeurs remarquables
			\foreach \a/\l/\pos in {%
				0/{1}/{below right},
				30/{\e^{\ii\pi/6}}/{right},
				45/{\e^{\ii\pi/4}}/{right},
				60/{\e^{\ii\pi/3}}/{above right},
				90/{\ii}/{above right},
				180/{-1}/{left},
				270/{-\ii}/{below}}
			{
				\filldraw[blue] (\a:1) circle (0.9pt)
				node[\pos, font=\scriptsize, blue, inner sep=2pt] {$\l$};
			}
		\end{tikzpicture}
	\end{center}
	\centerline{\footnotesize $\e^{\ii\theta}$ parcourt le cercle de centre $O$ et de rayon $1$.}
\end{resume}

\begin{theoreme}[Formule de Moivre]
	Pour tout $\theta \in \R$ et tout $n \in \Z$ :
	$$
	(\cos\theta + \ii\sin\theta)^n = \cos(n\theta) + \ii\sin(n\theta).
	$$
\end{theoreme}

\begin{exemple}
	Développer $\cos(2\theta)$ et $\sin(2\theta)$ à l'aide de la formule de Moivre.

	$(\cos\theta + \ii\sin\theta)^2 = \cos^2\theta - \sin^2\theta + 2\ii\cos\theta\sin\theta$,
	donc en identifiant parties réelle et imaginaire :
	$$
	\cos(2\theta) = \cos^2\theta - \sin^2\theta, \qquad \sin(2\theta) = 2\cos\theta\sin\theta.
	$$
\end{exemple}

% ------------------------------------------------------------
\begin{resume}{Formule de Moivre}
	\label{res:cplx_moivre}
	Pour tout $\theta \in \R$ et tout $n \in \Z$ :
	$$
	\boxed{(\cos\theta + \ii\sin\theta)^{n} = \cos(n\theta) + \ii\sin(n\theta)}
	\qquad\text{c'est-à-dire}\qquad
	\bigl(\e^{\ii\theta}\bigr)^{n} = \e^{\ii n\theta}.
	$$
	\textbf{À quoi ça sert :} \emph{développer} $\cos(n\theta)$ et $\sin(n\theta)$
	en puissances de $\cos\theta$ et $\sin\theta$ --- c'est l'opération inverse de
	la linéarisation.

	\smallskip
	\textbf{En 3 gestes.} \textbf{1.} écrire $(\cos\theta+\ii\sin\theta)^{n}$ ;
	\textbf{2.} développer avec le binôme de Newton ;
	\textbf{3.} identifier partie réelle et partie imaginaire.

	\smallskip
	\textbf{Exemple express ($n = 3$).}
	$\cos 3\theta = 4\cos^{3}\theta - 3\cos\theta$,
	\quad $\sin 3\theta = 3\sin\theta - 4\sin^{3}\theta$.

	\smallskip
	\begin{center}
		\begin{tikzpicture}[>=Stealth, scale=1.45]
			\draw[->] (-1.35,0) -- (1.5,0) node[right, font=\scriptsize] {$\Re$};
			\draw[->] (0,-1.3) -- (0,1.6) node[above, font=\scriptsize] {$\Im$};
			\draw[black!65] (0,0) circle (1);
			% puissances successives de z = e^{i pi/6}
			\foreach \a/\l/\pos in {%
				30/{z}/{right},
				60/{z^{2}}/{above right},
				90/{z^{3}}/{above right},
				120/{z^{4}}/{above left}}
			{
				\draw[blue!70, thin] (0,0) -- (\a:1);
				\filldraw[blue] (\a:1) circle (0.9pt)
				node[\pos, font=\scriptsize, blue, inner sep=2pt] {$\l$};
			}
			% arcs de rotation successifs (à l'intérieur du cercle)
			\foreach \d in {30,60,90}{
				\draw[->, red, thick]
				(\d+4:0.52) arc[start angle=\d+4, end angle=\d+26, radius=0.52];
			}
			\node[red, font=\scriptsize] at (45:0.70) {$\theta$};
		\end{tikzpicture}
	\end{center}
	\centerline{\footnotesize Avec $z = \e^{\ii\theta}$ : multiplier par $z$, c'est tourner de $\theta$.}
\end{resume}

\begin{propriete}{Formules d'Euler}{}
	Pour tout $\theta \in \R$ :
	$$
	\cos\theta = \frac{\e^{\ii\theta} + \e^{-\ii\theta}}{2},
	\qquad
	\sin\theta = \frac{\e^{\ii\theta} - \e^{-\ii\theta}}{2\ii}.
	$$
\end{propriete}

% ------------------------------------------------------------
\begin{resume}{Linéarisation}
	\label{res:cplx_linear}
	\textbf{But :} transformer un \emph{produit de puissances}
	$\cos^{p}\theta\,\sin^{q}\theta$ en une \emph{somme} de $\cos(k\theta)$ et
	$\sin(k\theta)$ --- indispensable pour \textbf{intégrer}.

	\smallskip
	\textbf{Outils.} $\cos\theta = \dfrac{\e^{\ii\theta}+\e^{-\ii\theta}}{2}$,
	\quad $\sin\theta = \dfrac{\e^{\ii\theta}-\e^{-\ii\theta}}{2\ii}$,
	\quad puis on regroupe :
	$$
	\e^{\ii k\theta} + \e^{-\ii k\theta} = 2\cos(k\theta),
	\qquad
	\e^{\ii k\theta} - \e^{-\ii k\theta} = 2\ii\,\sin(k\theta).
	$$

	\begin{center}
		\begin{tikzpicture}[>=Stealth, node distance=3.2mm,
			bloc/.style={draw=black!70, rounded corners=1mm, fill=gray!10,
				align=center, font=\scriptsize, inner sep=2pt, text width=1.68cm,
				minimum height=1cm}]
			\node[bloc] (a) {$\cos^{p}\theta\,\sin^{q}\theta$};
			\node[bloc, right=of a] (b) {\textbf{1.} Euler\\(remplacer)};
			\node[bloc, right=of b] (c) {\textbf{2.} Binôme\\de Newton};
			\node[bloc, right=of c] (d) {\textbf{3.} Regrouper\\les conjugués};
			\node[bloc, right=of d] (e) {$\sum \lambda_{k}\cos k\theta$};
			\draw[->, thick] (a) -- (b);
			\draw[->, thick] (b) -- (c);
			\draw[->, thick] (c) -- (d);
			\draw[->, thick] (d) -- (e);
		\end{tikzpicture}
	\end{center}

	\textbf{À connaître par cœur.}
	$$
	\cos^{2}\theta = \frac{1+\cos 2\theta}{2},
	\qquad
	\sin^{2}\theta = \frac{1-\cos 2\theta}{2},
	$$
	$$
	\cos^{3}\theta = \frac{\cos 3\theta + 3\cos\theta}{4},
	\qquad
	\sin^{3}\theta = \frac{3\sin\theta - \sin 3\theta}{4}.
	$$
	\textbf{Piège.} Avec le sinus, ne pas oublier le $\ii$ du dénominateur :
	$(2\ii)^{2} = -4$ et $(2\ii)^{3} = -8\ii$ --- c'est là que naissent les erreurs de signe.
\end{resume}

\begin{astucebac}
	\textbf{Applications des formules d'Euler et de Moivre.}
	\begin{itemize}[label=\textbullet]
		\item \textbf{Linéariser} $\cos^n\theta$, $\sin^n\theta$ : remplacer
		$\cos\theta$, $\sin\theta$ par les formules d'Euler, développer au
		binôme, regrouper les termes conjugués.
		\item \textbf{Développer} $\cos(n\theta)$, $\sin(n\theta)$ en fonction
		de $\cos\theta$, $\sin\theta$ : partir de la formule de Moivre
		$(\cos\theta + \ii\sin\theta)^n$ et identifier parties réelle et
		imaginaire.
		\item \textbf{Transformer un produit en somme} :
		$2\cos a\cos b = \cos(a-b) + \cos(a+b)$, etc.
	\end{itemize}
\end{astucebac}

\begin{propriete}{Réduction de $a\cos x + b\sin x$ sous la forme $R\cos(x-\varphi)$}{cplx_reduction}
	Pour $(a,b) \neq (0,0)$, il existe $R > 0$ et $\varphi \in \R$ tels que,
	pour tout $x$,
	$$
	a\cos x + b\sin x = R\cos(x - \varphi),
	$$
	avec $R = \sqrt{a^2 + b^2}$, $\cos\varphi = \dfrac{a}{R}$,
	$\sin\varphi = \dfrac{b}{R}$.

	\emph{Preuve :} $a\cos x + b\sin x = \mathrm{Re}\bigl((a - b\ii)\,\e^{\ii x}\bigr)$
	et $a - b\ii = R\,\e^{-\ii\varphi}$.
\end{propriete}


% ============================================================
\section{Équations polynomiales dans $\C$}
% ============================================================

\subsection{Racine carrée d'un nombre complexe}

\begin{methode}
	Pour déterminer les racines carrées de $Z = \alpha + \beta\ii$, on cherche
	$z = x + y\ii$ tel que $z^2 = Z$ en résolvant
	$$
	\begin{cases}
		x^2 - y^2 = \alpha & \text{(partie réelle)}\\
		2xy = \beta & \text{(partie imaginaire)}\\
		x^2 + y^2 = \sqrt{\alpha^2 + \beta^2} & \text{(modules)}
	\end{cases}
	$$
	Les deux premières équations donnent $x^2$ et $y^2$ ; le signe de $\beta$
	fixe le signe de $xy$.
\end{methode}

\begin{exemple}
	Racines carrées de $Z = 3 + 4\ii$ : $x^2 - y^2 = 3$, $2xy = 4 > 0$,
	$x^2 + y^2 = 5$. D'où $x^2 = 4$, $y^2 = 1$, $xy > 0$ :
	$z = 2 + \ii$ ou $z = -2 - \ii$.
\end{exemple}

\subsection{Équations du second degré}

\begin{theoreme}
	Soit $(E) : az^2 + bz + c = 0$ avec $a, b, c \in \C$, $a \neq 0$, et
	$\Delta = b^2 - 4ac$. On note $\delta$ \emph{une} racine carrée de $\Delta$
	dans $\C$.
	\begin{itemize}[label=\textbullet]
		\item Si $\Delta \neq 0$ : deux solutions
		$z = \dfrac{-b + \delta}{2a}$ et $z = \dfrac{-b - \delta}{2a}$.
		\item Si $\Delta = 0$ : une solution double $z = -\dfrac{b}{2a}$.
	\end{itemize}
	Dans tous les cas $z_1 + z_2 = -\dfrac{b}{a}$ et $z_1 z_2 = \dfrac{c}{a}$
	(relations coefficients–racines), et $az^2 + bz + c = a(z - z_1)(z - z_2)$.

	\textbf{Cas des coefficients réels} : si $a,b,c \in \R$ et $\Delta < 0$,
	les solutions sont \textbf{conjuguées} :
	$z_1 = \dfrac{-b + \ii\sqrt{-\Delta}}{2a}$, $z_2 = \overline{z_1}$.
\end{theoreme}

\begin{exemple}
	$z^2 + 2z + 5 = 0$ : $\Delta = -16$, $\delta = 4\ii$,
	$z_1 = -1 + 2\ii$, $z_2 = \overline{z_1} = -1 - 2\ii$.
\end{exemple}

\subsection{Équations de degré $3$}

\begin{methode}
	Pour résoudre $P(z) = 0$ avec $P$ de degré $3$ :
	\begin{enumerate}
		\item \textbf{Trouver une racine} $z_0$ : soit elle est donnée par
		l'énoncé (« vérifier que $z_0$ est solution »), soit on cherche une
		racine évidente (réelle : $0$, $\pm1$, $\pm2$… ; ou imaginaire pure).
		\item \textbf{Factoriser} : $P(z) = (z - z_0)\,Q(z)$ où $Q$ est de
		degré $2$, obtenu par identification ou division euclidienne :
		$$
		P(z) = (z - z_0)(a z^2 + \beta z + \gamma).
		$$
		\item \textbf{Résoudre} $Q(z) = 0$ (second degré).
	\end{enumerate}
\end{methode}

\begin{exemple}
	$P(z) = z^3 - 3z^2 + 4z - 2$. $P(1) = 1 - 3 + 4 - 2 = 0$, donc
	$P(z) = (z - 1)(z^2 - 2z + 2)$. Le facteur $z^2 - 2z + 2$ a pour
	discriminant $\Delta = -4$, racines $1 \pm \ii$.
	Solutions : $\{\,1\,;\,1+\ii\,;\,1-\ii\,\}$.
\end{exemple}

\subsection{Équations de degré $4$}

\begin{methode}
	\begin{itemize}[label=\textbullet]
		\item \textbf{Équation bicarrée} $z^4 + p z^2 + q = 0$ : poser
		$Z = z^2$, résoudre $Z^2 + pZ + q = 0$, puis extraire les racines
		carrées de chaque valeur de $Z$.
		\item \textbf{Deux racines connues} (souvent conjuguées $z_0$,
		$\overline{z_0}$) : $P(z) = (z - z_0)(z - \overline{z_0})\,(z^2 + \beta z + \gamma)$,
		on identifie $\beta$, $\gamma$.
		\item \textbf{Factorisation en deux trinômes réels} : chercher
		$z^4 + \cdots = (z^2 + \alpha z + \beta)(z^2 + \alpha' z + \beta')$
		par identification.
	\end{itemize}
\end{methode}

\begin{exemple}
	$z^4 + 4 = 0 \iff z^4 = -4 = 4\,\e^{\ii\pi}$. Les solutions sont
	$z_k = \sqrt{2}\,\e^{\ii\left(\frac{\pi}{4} + \frac{k\pi}{2}\right)}$,
	$k \in \{0,1,2,3\}$, soit $\pm(1 + \ii)$ et $\pm(1 - \ii)$.
\end{exemple}

\subsection{Racines $n$-ièmes}

\begin{theoreme}[Racines $n$-ièmes de l'unité]
	Les solutions de $z^n = 1$ sont les $n$ nombres
	$$
	\omega_k = \e^{\ii\frac{2k\pi}{n}}, \quad k \in \{0, 1, \ldots, n-1\}.
	$$
	Leurs images sont les sommets d'un polygone régulier à $n$ côtés inscrit
	dans le cercle unité, et $\displaystyle\sum_{k=0}^{n-1}\omega_k = 0$.
\end{theoreme}

\begin{propriete}{Racines $n$-ièmes d'un complexe non nul}{cplx_racines_n}
	Soit $Z = \rho\,\e^{\ii\varphi}$ avec $\rho > 0$. Les solutions de $z^n = Z$
	sont les $n$ nombres
	$$
	z_k = \rho^{1/n}\,\e^{\ii\left(\frac{\varphi}{n} + \frac{2k\pi}{n}\right)},
	\quad k \in \{0,1,\ldots,n-1\}.
	$$
	On les obtient à partir d'une solution particulière en multipliant
	successivement par les racines $n$-ièmes de l'unité.
\end{propriete}

\begin{exemple}
	Les racines cubiques de l'unité sont
	$1$, $\;j = \e^{\ii\frac{2\pi}{3}} = -\dfrac{1}{2} + \dfrac{\sqrt3}{2}\ii$,
	$\;j^2 = \overline{j}$, avec $1 + j + j^2 = 0$.
\end{exemple}


% ============================================================
%\section{Systèmes de deux équations à deux inconnues}
%% ============================================================
%
%\subsection{Système somme–produit}
%
%\begin{theoreme}
%	Deux complexes $z$ et $z'$ vérifient
%	$$
%	\begin{cases}
%		z + z' = S \\
%		z\,z' = P
%	\end{cases}
%	$$
%	si et seulement si $z$ et $z'$ sont les deux solutions (éventuellement
%	confondues) de l'équation
%	$$
%	X^2 - S\,X + P = 0.
%	$$
%\end{theoreme}
%
%\begin{exemple}
%	$\begin{cases} z + z' = 4 \\ z\,z' = 5 \end{cases}$ : $X^2 - 4X + 5 = 0$,
%	$\Delta = -4$, donc $\{z, z'\} = \{\,2 + \ii\,;\,2 - \ii\,\}$.
%\end{exemple}
%
%\subsection{Système linéaire à coefficients complexes}
%
%\begin{methode}
%	Pour $\begin{cases} a z + b z' = c \\ a' z + b' z' = c' \end{cases}$
%	($a,b,\ldots \in \C$) : on résout comme dans $\R$ (combinaisons linéaires,
%	substitution), en veillant à ne diviser que par un complexe non nul.
%
%	\textbf{Systèmes faisant intervenir $\bar z$} : quand une équation porte
%	sur $z$ et $\bar z$, on prend le \textbf{conjugué} de l'équation pour
%	obtenir un second lien, puis on résout le système en $(z, \bar z)$.
%\end{methode}
%
%\begin{exemple}
%	$\begin{cases} z + \ii z' = 1 + \ii \\ \ii z + z' = 2 \end{cases}$.
%	En multipliant la première ligne par $\ii$ : $\ii z - z' = \ii - 1$.
%	En soustrayant la seconde : $-2z' = -1 + \ii - 2 = -3 + \ii$, d'où
%	$z' = \dfrac{3 - \ii}{2}$, puis $z = 2 - z' = \dfrac{1 + \ii}{2}$.
%\end{exemple}
%
%
%% ============================================================
%\section{Nombres complexes et géométrie}
% ============================================================

\begin{propriete}{Affixe d'un point et d'un vecteur}{}
	Dans le plan complexe repéré par $(O, \vec{\imath}, \vec{\jmath})$ :
	\begin{itemize}[label=\textbullet]
		\item Le point $M$ d'affixe $z = a + b\ii$ a pour coordonnées $(a\,;\,b)$.
		\item Le vecteur $\vect{AB}$ a pour affixe $z_B - z_A$.
		\item La distance $AB = \abs{z_B - z_A}$.
		\item Le milieu $I$ de $[AB]$ a pour affixe $\dfrac{z_A + z_B}{2}$.
	\end{itemize}
\end{propriete}


\exerciceresolu
	\begin{figure}[h]
		\centering
		\begin{tikzpicture}[>=Stealth, scale=1.1]

			% --- Grille -----------------------------------------------
			\draw[gray!30, thin] (0,0) grid (6,3);

			% --- Axes -------------------------------------------------
			\draw[->] (-0.3,0) -- (6.5,0) node[right, font=\small] {Axe des réels};
			\draw[->] (0,-0.3) -- (0,3.5) node[above, font=\small] {Axe des imaginaires};

			% --- Graduation axe réel ----------------------------------
			\foreach \x in {1,...,6}{
				\draw (\x, 2pt) -- (\x, -2pt) node[below, font=\small] {\x};
			}
			% zéro
			\node[below left, font=\small] at (0,0) {$0$};

			% --- Graduation axe imaginaire ----------------------------
			\foreach \y in {1,2,3}{
				\draw (2pt,\y) -- (-2pt,\y) node[left, font=\small] {\y};
			}

			% --- Vecteurs de base u et v (en rouge) -------------------
			\draw[->, red, thick] (0,0) -- (1,0)
			node[below, midway, font=\small, red] {$\vec{u}$};
			\draw[->, red, thick] (0,0) -- (0,1)
			node[midway, left,  font=\small, red] {$\vec{v}$};

			% --- Tirets pointillés vers M(3,2) -----------------------
			\draw[black!60, dashed, thin] (3,0) -- (3,2);
			\draw[black!60, dashed, thin] (0,2) -- (3,2);

			% --- Vecteur w = affixe 3+2i (en vert) -------------------
			\draw[->, blue!60!black, thick] (0,0) -- (3,2)
			node[midway, below right, font=\small, blue!60!black]
			{$\vec{w}(3+2i)$};



			% --- Point M ----------------------------------------------
			\filldraw[blue] (3,2) circle (2pt)
			node[above right, font=\small, blue] {$M(3+2i)$};

		\end{tikzpicture}

		\smallskip
		\begin{minipage}{0.85\linewidth}
			\small
			Le point $M\!\begin{pmatrix}3\\2\end{pmatrix}$ a pour affixe le nombre complexe
			$z = 3 + 2i$.\\
			De même, le vecteur $\vec{w}$ a pour affixe $z = 3 + 2i$.\\
			La longueur $OM$ est {$OM = |w| =\sqrt{3^2+2^2}=\sqrt{13}$}
		\end{minipage}
%		\caption{Affixe d'un point et d'un vecteur dans le plan complexe.}
%		\label{fig:affixe}
	\end{figure}
%\end{exemple}

\begin{propriete}{Le quotient $\dfrac{z_C - z_A}{z_B - z_A}$}{cplx_quotient}
	Soient $A$, $B$, $C$ trois points distincts. En posant
	$Z = \dfrac{z_C - z_A}{z_B - z_A}$ :
	$$
	\abs{Z} = \frac{AC}{AB},
	\qquad
	\arg(Z) \equiv \bigl(\vect{AB},\,\vect{AC}\bigr)\ [2\pi].
	$$
\end{propriete}

\begin{propriete}{Conditions géométriques}{cplx_conditions}
	Soient $A$, $B$, $C$, $D$ des points deux à deux distincts.
	\begin{enumerate}
		\item $A$, $B$, $C$ \textbf{alignés} $\iff \dfrac{z_C - z_A}{z_B - z_A} \in \R$.
		\item $(AB) \perp (CD) \iff \dfrac{z_D - z_C}{z_B - z_A} \in \ii\R$
		(imaginaire pur).
		\item $(AB) \parallel (CD) \iff \dfrac{z_D - z_C}{z_B - z_A} \in \R$.
		\item $ABDC$ \textbf{parallélogramme} $\iff z_D - z_C = z_B - z_A$.
		\item $A$, $B$, $C$, $D$ \textbf{cocycliques} (ou alignés) $\iff$
		$\dfrac{z_C - z_A}{z_D - z_A} \times \dfrac{z_D - z_B}{z_C - z_B} \in \R$.
	\end{enumerate}
\end{propriete}

\begin{propriete}{Nature d'un triangle $ABC$}{cplx_triangle}
	En posant $Z = \dfrac{z_C - z_A}{z_B - z_A}$ :
	\begin{itemize}[label=\textbullet]
		\item $ABC$ \textbf{isocèle en $A$} $\iff \abs{Z} = 1$ ;
		\item $ABC$ \textbf{rectangle en $A$} $\iff Z \in \ii\R$ ;
		\item $ABC$ \textbf{rectangle isocèle en $A$} $\iff Z = \pm\ii$ ;
		\item $ABC$ \textbf{équilatéral} $\iff Z = \e^{\pm\ii\pi/3}$, ce qui
		équivaut à $z_A^2 + z_B^2 + z_C^2 = z_A z_B + z_B z_C + z_C z_A$.
	\end{itemize}
\end{propriete}

\begin{exemple}
	$A(1+\ii)$, $B(3+2\ii)$ : milieu $I$ d'affixe $\dfrac{4 + 3\ii}{2}$ ;
	$AB = \abs{2 + \ii} = \sqrt{5}$.
\end{exemple}


% ============================================================
\section{Ensembles de points}
% ============================================================

Le plan complexe est rapporté à un repère orthonormé direct
$(O,\vec u,\vec v)$. On cherche l'ensemble des points $M$ d'affixe $z$
vérifiant une condition donnée. On note $A$, $B$ les points d'affixes $a$, $b$.

\subsection{Conditions de module (distances)}

\begin{propriete}{}{cplx_ens_module}
	\begin{enumerate}
		\item $\abs{z - a} = r$ \; ($r > 0$) : \textbf{cercle} de centre $A$,
		de rayon $r$. \\ ($\abs{z-a} \leqslant r$ : disque fermé.)
		\item $\abs{z - a} = \abs{z - b}$ : \textbf{médiatrice} de $[AB]$.
		\item $\abs{z - a} = k\,\abs{z - b}$ \; ($k > 0$, $k \neq 1$) :
		\textbf{cercle} (dit d'Apollonius) ; si $k = 1$ : la médiatrice de $[AB]$.
		\item $\abs{z} = \abs{z - a}$, plus généralement toute égalité de deux
		modules, se ramène à l'un des cas ci-dessus.
	\end{enumerate}
\end{propriete}

\begin{center}
\begin{tikzpicture}[>=Stealth, scale=0.95]
	% cercle |z-a| = r
	\begin{scope}
		\draw[->] (-1.4,0)--(2.2,0); \draw[->] (0,-1.4)--(0,1.7);
		\filldraw (0.8,0.4) circle (1.3pt) node[above right,font=\scriptsize] {$A$};
		\draw[blue!70!black,thick] (0.8,0.4) circle (1);
		\node[font=\scriptsize] at (0.4,-1.7) {$\abs{z-a}=r$};
	\end{scope}
	% médiatrice
	\begin{scope}[xshift=4.2cm]
		\draw[->] (-1.4,0)--(2.2,0); \draw[->] (0,-1.4)--(0,1.7);
		\filldraw (-0.7,-0.4) circle (1.3pt) node[below left,font=\scriptsize] {$A$};
		\filldraw (1.1,0.6) circle (1.3pt) node[above right,font=\scriptsize] {$B$};
		\draw[gray!60,dashed] (-0.7,-0.4)--(1.1,0.6);
		\draw[blue!70!black,thick] (-0.6,1.7) -- (1.0,-1.5);
		\node[font=\scriptsize] at (0.2,-1.9) {$\abs{z-a}=\abs{z-b}$};
	\end{scope}
\end{tikzpicture}
\end{center}

\subsection{Conditions d'argument}

\begin{propriete}{}{cplx_ens_arg}
	\begin{enumerate}
		\item $\arg(z - a) \equiv \theta\ [2\pi]$ : \textbf{demi-droite}
		d'origine $A$ (exclue), dirigée par l'angle $\theta$. \\
		($\arg(z-a) \equiv \theta\ [\pi]$ : la \textbf{droite} passant par $A$
		privée de $A$.)
		\item $\arg\!\left(\dfrac{z - a}{z - b}\right) \equiv \theta\ [2\pi]$
		(avec $\theta \not\equiv 0\ [\pi]$) : \textbf{arc de cercle} passant
		par $A$ et $B$ (extrémités exclues).
		\item $\dfrac{z - a}{z - b} \in \R^*$ : la \textbf{droite} $(AB)$
		privée de $A$ et $B$.
		\item $\dfrac{z - a}{z - b} \in \ii\R^*$ : le \textbf{cercle} de
		diamètre $[AB]$ privé de $A$ et $B$.
	\end{enumerate}
\end{propriete}

\begin{center}
\begin{tikzpicture}[>=Stealth, scale=0.95]
	% demi-droite arg(z-a) = theta
	\begin{scope}
		\draw[->] (-1.4,0)--(2.2,0); \draw[->] (0,-1.4)--(0,1.7);
		\filldraw (-0.4,-0.3) circle (1.3pt) node[below,font=\scriptsize] {$A$};
		\draw[blue!70!black,thick,->] (-0.4,-0.3) -- (1.8,1.35);
		\draw[red] (-0.4,-0.3)++(0:0.55) arc (0:36:0.55);
		\node[red,font=\scriptsize] at (0.45,-0.15) {$\theta$};
		\node[font=\scriptsize] at (0.4,-1.75) {$\arg(z-a)=\theta$};
	\end{scope}
	% cercle de diamètre [AB]
	\begin{scope}[xshift=4.2cm]
		\draw[->] (-1.6,0)--(2.0,0); \draw[->] (0,-1.5)--(0,1.6);
		\filldraw (-0.9,-0.2) circle (1.3pt) node[left,font=\scriptsize] {$A$};
		\filldraw (0.9,0.4) circle (1.3pt) node[right,font=\scriptsize] {$B$};
		\draw[blue!70!black,thick] (0,0.1) circle ({sqrt(0.9*0.9+0.3*0.3)});
		\draw[gray!60,dashed] (-0.9,-0.2)--(0.9,0.4);
		\node[font=\scriptsize] at (0.1,-1.9) {$\frac{z-a}{z-b}\in\ii\R$};
	\end{scope}
\end{tikzpicture}
\end{center}

\subsection{Conditions sur la partie réelle ou imaginaire}

\begin{propriete}{}{cplx_ens_reim}
	\begin{enumerate}
		\item $\mathrm{Re}(z) = k$ : droite \textbf{verticale} $x = k$ ;
		\quad $\mathrm{Im}(z) = k$ : droite \textbf{horizontale} $y = k$.
		\item $z + \bar z = k \iff \mathrm{Re}(z) = \frac{k}{2}$ ;
		\quad $z - \bar z = k\ii \iff \mathrm{Im}(z) = \frac{k}{2}$.
		\item $\mathrm{Re}\!\left(\dfrac{z - a}{z - b}\right) = 0$ : cercle de
		diamètre $[AB]$ (privé de $A$, $B$) ;
		$\;\mathrm{Im}\!\left(\dfrac{z - a}{z - b}\right) = 0$ : droite $(AB)$
		(privée de $A$, $B$).
		\item $z\bar z + \bar\alpha z + \alpha\bar z + \beta = 0$
		($\beta \in \R$) : cercle (ou point, ou $\emptyset$), car cela s'écrit
		$\abs{z + \alpha}^2 = \abs{\alpha}^2 - \beta$.
	\end{enumerate}
\end{propriete}

\begin{astucebac}
	\textbf{Méthode générale.}
	\begin{enumerate}
		\item Traduire la condition avec $z = x + \ii y$ \emph{ou} l'interpréter
		directement (distance $\abs{z - a} = AM$, angle
		$\arg\frac{z-c}{z-d} = (\vect{DC},\ldots)$).
		\item Reconnaître : distance constante $\to$ cercle ; distances égales
		$\to$ médiatrice ; angle droit $\to$ cercle de diamètre ; angle
		constant $\to$ arc ; quotient réel $\to$ droite.
		\item Préciser les \textbf{éléments} (centre, rayon, points exclus) et
		\textbf{construire}.
	\end{enumerate}
\end{astucebac}

\begin{center}
\renewcommand{\arraystretch}{1.35}
\begin{tabular}{>{\raggedright\arraybackslash}p{4.0cm}
                >{\raggedright\arraybackslash}p{4.2cm}}
\hline
\textbf{Condition} & \textbf{Ensemble} \\
\hline
$\abs{z - a} = r$ & cercle $C(A, r)$ \\
$\abs{z - a} = \abs{z - b}$ & médiatrice de $[AB]$ \\
$\abs{z - a} = k\abs{z - b}$, $k \neq 1$ & cercle (Apollonius) \\
$\arg(z - a) = \theta\ [2\pi]$ & demi-droite d'origine $A$ \\
$\dfrac{z - a}{z - b} \in \R$ & droite $(AB)$ privée de $A$, $B$ \\
$\dfrac{z - a}{z - b} \in \ii\R$ & cercle de diamètre $[AB]$ \\
$\mathrm{Re}(z) = k$ / $\mathrm{Im}(z) = k$ & droite verticale / horizontale \\
\hline
\end{tabular}
\end{center}


% ============================================================
\section{Transformations du plan et complexes}
% ============================================================

\begin{propriete}{Translations}{}
	La translation de vecteur $\vect{u}$ d'affixe $b$ est donnée par :
	$$
	z' = z + b.
	$$
\end{propriete}

\begin{propriete}{Rotation}{}
	La rotation de centre $\Omega$ (affixe $\omega$), d'angle $\theta$ est donnée par :
	$$
	z' - \omega = \e^{\ii\theta}(z - \omega).
	$$
\end{propriete}

\begin{propriete}{Homothétie}{}
	L'homothétie de centre $\Omega$ (affixe $\omega$), de rapport $k \in \R^*$ est :
	$$
	z' - \omega = k(z - \omega).
	$$
\end{propriete}

\begin{propriete}{Similitude directe}{}
	Toute similitude directe (non nulle) s'écrit sous la forme
	$$
	z' = az + b, \quad a \in \C^*,\ b \in \C.
	$$
	\begin{itemize}[label=\textbullet]
		\item Le \textbf{rapport} de la similitude est $\abs{a}$.
		\item L'\textbf{angle} de la similitude est $\arg(a)$.
		\item Si $a = 1$ : c'est une translation.
		\item Si $a \neq 1$ : le \textbf{centre} est le point fixe $\omega = \dfrac{b}{1 - a}$.
	\end{itemize}
\end{propriete}

\begin{exemple}
	La transformation $f : z \mapsto (1 + \ii)z - 2 + \ii$ est une similitude directe.
	\begin{itemize}[label=\textbullet]
		\item Rapport : $\abs{1 + \ii} = \sqrt{2}$.
		\item Angle : $\arg(1 + \ii) = \dfrac{\pi}{4}$.
		\item Centre : $\omega = \dfrac{-2 + \ii}{1 - (1 + \ii)} = \dfrac{-2 + \ii}{-\ii}
		= \dfrac{(-2 + \ii)\ii}{-\ii^2} = \dfrac{-2\ii + \ii^2}{1} = -1 - 2\ii$.
	\end{itemize}
\end{exemple}

%\begin{astucebac}
%	Pour identifier une transformation à partir de son expression complexe $z' = az + b$ :
%	\begin{enumerate}
%		\item Calculer $a$ : son module donne le rapport, son argument donne l'angle.
%		\item Si $a = 1$ : translation de vecteur d'affixe $b$.
%		\item Si $\abs{a} = 1$ et $a \neq 1$ : rotation pure (rapport $= 1$).
%		\item Si $a \in \R^*$ et $a \neq 1$ : homothétie pure (angle $= 0$ ou $\pi$).
%		\item Sinon : similitude directe de rapport $\abs{a}$ et d'angle $\arg(a)$.
%	\end{enumerate}
%\end{astucebac}

\begin{astucebac}
	Pour identifier une transformation à partir de son expression complexe $z' = az + b$ :

	\begin{center}
		\begin{tikzpicture}[
			>=Latex,
			font=\small,
			start/.style={draw, rounded corners, minimum width=3.6cm, minimum height=0.8cm, align=center},
			decision/.style={draw, diamond, aspect=2.2, minimum width=3.4cm, minimum height=1cm, align=center, inner sep=2pt},
			resultat/.style={draw, rounded corners, minimum width=3.8cm, minimum height=0.9cm, align=center},
			lab/.style={font=\scriptsize}
			]

			% --- Etape de départ ---
			\node[start] (depart) {Calculer $a$\\(module et argument)};

			% --- Décision 1 ---
			\node[decision, below=0.9cm of depart] (d1) {$a = 1$ ?};
			\node[resultat, right=1.6cm of d1] (r1) {Translation\\de vecteur d'affixe $b$};

			% --- Décision 2 ---
			\node[decision, below=1.3cm of d1] (d2) {$\abs{a}=1$ et $a\neq 1$ ?};
			\node[resultat, right=1.6cm of d2] (r2) {Rotation pure\\(rapport $=1$, angle $\arg(a)$)};

			% --- Décision 3 ---
			\node[decision, below=1.3cm of d2] (d3) {$a\in\R^*$ et $a\neq 1$ ?};
			\node[resultat, right=1.6cm of d3] (r3) {Homothétie pure\\(angle $=0$ ou $\pi$)};

			% --- Résultat final ---
			\node[resultat, below=1.3cm of d3] (r4) {Similitude directe\\de rapport $\abs{a}$ et d'angle $\arg(a)$};

			% --- flèches verticales (chemin "non") ---
			\draw[->] (depart) -- (d1);
			\draw[->] (d1) -- node[lab, left]{non} (d2);
			\draw[->] (d2) -- node[lab, left]{non} (d3);
			\draw[->] (d3) -- node[lab, left]{non} (r4);

			% --- flèches horizontales (chemin "oui") ---
			\draw[->] (d1) -- node[lab, above]{oui} (r1);
			\draw[->] (d2) -- node[lab, above]{oui} (r2);
			\draw[->] (d3) -- node[lab, above]{oui} (r3);

		\end{tikzpicture}
	\end{center}
\end{astucebac}



% ============================================================
%   EXERCICES RÉSOLUS — Nombres complexes
%   Mazava Maths
% ============================================================


% ------------------------------------------------------------
%  NOTION 1 : Forme algébrique — opérations et conjugué
% ------------------------------------------------------------



\section{Méthodes et exercices résolus}

\begin{methode}
	\textbf{Résoudre une équation polynomiale de degré $\geqslant 3$.}
	\begin{enumerate}
		\item Utiliser (ou chercher) une racine $z_0$.
		\item Factoriser $P(z) = (z - z_0)\,Q(z)$ par identification ou en utilisant le schéma d'Horner .
		\item Résoudre $Q(z) = 0$. Pour un degré $4$ : deux racines connues, ou
		poser $Z = z^2$ si l'équation est bicarrée.
	\end{enumerate}
\end{methode}

\begin{methode}
	\textbf{Déterminer un ensemble de points.}
	Interpréter la condition en termes de \emph{distance} ($\abs{z - a} = AM$)
	ou d'\emph{angle} ($\arg\frac{z - a}{z - b}$) ; sinon poser $z = x + \ii y$
	et traduire. Reconnaître : cercle, médiatrice, droite, demi-droite, arc,
	cercle de diamètre. Préciser les éléments et \textbf{construire}.
\end{methode}

\exerciceresolu
\textbf{Équation de degré $3$}

Résoudre dans $\C$ : $z^3 - (4 + \ii)z^2 + (5 + 4\ii)z - 5\ii = 0$, sachant
qu'elle admet une racine imaginaire pure.

\begin{solution}
	Soit $z = \ii \alpha$ ($\alpha \in \R$) une racine candidate. En reportant :
	$$
	-\ii \alpha^3 + (4 + \ii)\alpha^2 + (5 + 4\ii)\ii \alpha - 5\ii = 0,
	$$
	soit après regroupement
	$$
	\bigl(4\alpha^2 - 4\alpha\bigr) + \ii\bigl(-\alpha^3 + \alpha^2 + 5\alpha - 5\bigr) = 0.
	$$
	La partie réelle donne $4\alpha(\alpha - 1) = 0$ ; $\alpha = 1$ annule aussi la
	partie imaginaire. Donc $z_0 = \ii$ est racine.

	\medskip
	\textbf{Factorisation par la méthode de Horner.} On divise
	$z^3 - (4+\ii)z^2 + (5+4\ii)z - 5\ii$ par $(z - \ii)$ :

	\begin{center}
		\begin{tikzpicture}[>=Latex, font=\small]
			% --- coefficients du polynôme ---
			\node (c0) at (0,0)   {$1$};
			\node (c1) at (2,0)   {$-(4+\ii)$};
			\node (c2) at (4.4,0) {$5+4\ii$};
			\node (c3) at (6.4,0) {$-5\ii$};

			% --- racine ---
			\node[left] at (-1.2,-1) {racine $\ii$};

			% --- ligne des produits (racine x quotient précédent) ---
			\node (p1) at (2,-1)   {$\ii$};
			\node (p2) at (4.4,-1) {$-4\ii$};
			\node (p3) at (6.4,-1) {$5\ii$};

			% --- ligne des résultats (coefficients du quotient + reste) ---
			\node (q0) at (0,-2)   {$1$};
			\node (q1) at (2,-2)   {$-4$};
			\node (q2) at (4.4,-2) {$5$};
			\node (q3) at (6.4,-2) {$0$};

			% --- traits horizontaux du schéma ---
			\draw (-0.6,-0.5) -- (7.2,-0.5);
			\draw (-0.6,-1.5) -- (7.2,-1.5);

			% --- flèches diagonales : "descendre" puis "multiplier par i" ---
			\draw[->] (c0) -- (p1);
			\draw[->] (c1) -- (p2);
			\draw[->] (c2) -- (p3);

			% --- flèches verticales : addition ---
			\draw[->] (c0) -- (q0);
			\draw[->] (p1) -- (q1);
			\draw[->] (p2) -- (q2);
			\draw[->] (p3) -- (q3);

		\end{tikzpicture}
	\end{center}

	Le reste est nul, ce qui confirme que $\ii$ est bien racine, et le quotient
	donne directement :
	$$
	z^3 - (4+\ii)z^2 + (5+4\ii)z - 5\ii = (z - \ii)(z^2 - 4z + 5).
	$$

	\medskip
	Il reste à résoudre $z^2 - 4z + 5 = 0$ : $\Delta = -4$, donc $z = 2 \pm \ii$.

	\medskip
	\textbf{Conclusion.} Les trois racines sont $z_0 = \ii$, $z_1 = 2+\ii$,
	$z_2 = 2-\ii$.
\end{solution}

\exerciceresolu
\textbf{Effectuer des calculs en forme algébrique}

Soient $z_1 = 3 - 2\ii$ et $z_2 = 1 + 4\ii$.

Calculer, sous forme algébrique, $z_1 + z_2$, $z_1 \times z_2$, $\overline{z_1}$
et $\dfrac{z_1}{z_2}$.

\solution

\textbf{Somme :}
$$
z_1 + z_2 = (3+1) + (-2+4)\ii = 4 + 2\ii.
$$

\textbf{Produit :}
$$
z_1 \times z_2 = (3 \times 1 - (-2) \times 4) + (3 \times 4 + (-2) \times 1)\ii
= (3 + 8) + (12 - 2)\ii = 11 + 10\ii.
$$

\textbf{Conjugué :}
$$
\overline{z_1} = 3 + 2\ii.
$$

\textbf{Quotient :} on multiplie par le conjugué du dénominateur $\overline{z_2} = 1 - 4\ii$ :
$$
\frac{z_1}{z_2}
= \frac{(3 - 2\ii)(1 - 4\ii)}{(1 + 4\ii)(1 - 4\ii)}
= \frac{3 - 12\ii - 2\ii + 8\ii^2}{1 + 16}
= \frac{3 - 8 - 14\ii}{17}
= \frac{-5}{17} - \frac{14}{17}\ii.
$$


% ------------------------------------------------------------
%  NOTION 2 : Résoudre une équation à coefficients réels
% ------------------------------------------------------------

\exerciceresolu
\textbf{Résoudre une équation et trouver des parties réelle et imaginaire}

Résoudre dans $\C$ l'équation $z + \ii\bar{z} = 3 - \ii$.

\solution

On pose $z = a + b\ii$ avec $a, b \in \R$, donc $\bar{z} = a - b\ii$.

On substitue :
$$
\begin{aligned}
	(a + b\ii) + \ii(a - b\ii) &= 3 - \ii\\
	a + b\ii + a\ii - b\ii^2   &= 3 - \ii\\
	(a + b) + (a + b)\ii       &= 3 - \ii
\end{aligned}
$$

En identifiant parties réelle et imaginaire :
$$
\begin{cases}
	a + b = 3 \\
	a + b = -1
\end{cases}
$$

Ce système est contradictoire ($3 \neq -1$), donc \textbf{l'équation n'a pas de solution}.


% ------------------------------------------------------------
%  NOTION 3 : Module
% ------------------------------------------------------------

\exerciceresolu
\textbf{Calculer un module et utiliser ses propriétés}

\begin{enumerate}
	\item Calculer $\abs{z}$ pour $z = \dfrac{(2 + \ii)^2}{1 - \ii\sqrt{3}}$.

	\item Montrer que si $\abs{z} = 1$, alors $\dfrac{z - 1}{z + 1}$ est imaginaire pur
	(en supposant $z \neq -1$).
\end{enumerate}

\solution

\textbf{1.} On utilise les propriétés du module :
$$
\abs{z}
= \frac{\abs{2 + \ii}^2}{\abs{1 - \ii\sqrt{3}}}
= \frac{\left(\sqrt{4+1}\right)^2}{\sqrt{1 + 3}}
= \frac{5}{2}.
$$

\textbf{2.} Posons $w = \dfrac{z-1}{z+1}$. Montrons que $\mathrm{Re}(w) = 0$.

On calcule $w + \bar{w}$ :
$$
w + \bar{w} = \frac{z-1}{z+1} + \frac{\bar{z}-1}{\bar{z}+1}
= \frac{(z-1)(\bar{z}+1) + (\bar{z}-1)(z+1)}{(z+1)(\bar{z}+1)}.
$$

Développons le numérateur :
$$
(z-1)(\bar{z}+1) + (\bar{z}-1)(z+1)
= z\bar{z} + z - \bar{z} - 1 + \bar{z}z + \bar{z} - z - 1
= 2\abs{z}^2 - 2.
$$

Comme $\abs{z} = 1$, on obtient $2 \times 1 - 2 = 0$.

Donc $w + \bar{w} = 0$, ce qui signifie $2\,\mathrm{Re}(w) = 0$, soit $\mathrm{Re}(w) = 0$.

Ainsi $\dfrac{z-1}{z+1}$ est bien imaginaire pur.


% ------------------------------------------------------------
%  NOTION 4 : Forme trigonométrique et argument
% ------------------------------------------------------------

\exerciceresolu
\textbf{Mettre sous forme trigonométrique et exponentielle}

Mettre les nombres complexes suivants sous forme exponentielle $r\e^{\ii\theta}$ :
\begin{listecol}(2)
	\task $z_1 = -1 + \ii$
	\task $z_2 = -2\ii$
	\task $z_3 = \sqrt{3} - \ii$
\end{listecol}

\solution

\textbf{1.} $\abs{z_1} = \sqrt{1+1} = \sqrt{2}$.
$$
z_1 = \sqrt{2}\left(-\frac{1}{\sqrt{2}} + \frac{1}{\sqrt{2}}\ii\right)
= \sqrt{2}\left(\cos\frac{3\pi}{4} + \ii\sin\frac{3\pi}{4}\right)
= \sqrt{2}\,\e^{\ii\frac{3\pi}{4}}.
$$

\textbf{2.} $\abs{z_2} = 2$.
$$
z_2 = 2\left(0 - \ii\right) = 2\left(\cos\left(-\frac{\pi}{2}\right)
+ \ii\sin\left(-\frac{\pi}{2}\right)\right) = 2\,\e^{-\ii\frac{\pi}{2}}.
$$

\textbf{3.} $\abs{z_3} = \sqrt{3+1} = 2$.
$$
z_3 = 2\left(\frac{\sqrt{3}}{2} - \frac{1}{2}\ii\right)
= 2\left(\cos\left(-\frac{\pi}{6}\right) + \ii\sin\left(-\frac{\pi}{6}\right)\right)
= 2\,\e^{-\ii\frac{\pi}{6}}.
$$


% ------------------------------------------------------------
%  NOTION 4 bis : Puissances avec la forme exponentielle
% ------------------------------------------------------------

\exerciceresolu
\textbf{Calculer une puissance avec la forme exponentielle}

Calculer $(1 + \ii)^{12}$ puis $\left(\sqrt{3} - \ii\right)^{9}$.

\solution

\emph{Principe} (résumé p.~\pageref{res:cplx_expo}) : on passe en forme
exponentielle, car $\bigl(r\e^{\ii\theta}\bigr)^{n} = r^{n}\e^{\ii n\theta}$.

\textbf{1.} $\abs{1+\ii} = \sqrt{2}$ et $\arg(1+\ii) = \dfrac{\pi}{4}$, donc
$1 + \ii = \sqrt{2}\,\e^{\ii\pi/4}$. Alors
$$
(1+\ii)^{12} = \left(\sqrt{2}\right)^{12}\e^{\ii\,12\times\frac{\pi}{4}}
= 2^{6}\,\e^{3\ii\pi} = 64\,\e^{\ii\pi} = -64.
$$

\textbf{2.} $\abs{\sqrt3 - \ii} = 2$ et $\arg(\sqrt3-\ii) = -\dfrac{\pi}{6}$,
donc $\sqrt3 - \ii = 2\,\e^{-\ii\pi/6}$. Alors
$$
\left(\sqrt3 - \ii\right)^{9} = 2^{9}\,\e^{-\ii\,9\times\frac{\pi}{6}}
= 512\,\e^{-\ii\frac{3\pi}{2}} = 512\,\e^{\ii\frac{\pi}{2}} = 512\,\ii.
$$

\begin{attention}
	Le calcul direct au binôme est ici hors de portée : c'est la forme
	exponentielle qui rend l'exercice faisable. Penser à ramener l'angle
	final dans $]-\pi\,;\,\pi]$ modulo $2\pi$ avant de conclure.
\end{attention}


% ------------------------------------------------------------
%  NOTION 5 : Formule de Moivre — linéarisation
% ------------------------------------------------------------

\exerciceresolu
\textbf{Exprimer $\cos(3\theta)$ et $\sin(3\theta)$ en fonction de $\cos\theta$ et $\sin\theta$}

\solution

D'après la formule de Moivre :
$$
(\cos\theta + \ii\sin\theta)^3 = \cos(3\theta) + \ii\sin(3\theta).
$$

On développe le membre gauche par le binôme de Newton :
$$
\begin{aligned}
	(\cos\theta + \ii\sin\theta)^3
	&= \cos^3\theta + 3\cos^2\theta\cdot\ii\sin\theta
	 + 3\cos\theta\cdot(\ii\sin\theta)^2 + (\ii\sin\theta)^3\\
	&= \cos^3\theta + 3\ii\cos^2\theta\sin\theta
	 - 3\cos\theta\sin^2\theta - \ii\sin^3\theta\\
	&= \left(\cos^3\theta - 3\cos\theta\sin^2\theta\right)
	 + \ii\left(3\cos^2\theta\sin\theta - \sin^3\theta\right)
\end{aligned}
$$

En identifiant parties réelle et imaginaire :
$$
\boxed{\cos(3\theta) = \cos^3\theta - 3\cos\theta\sin^2\theta = 4\cos^3\theta - 3\cos\theta,}
$$
$$
\boxed{\sin(3\theta) = 3\cos^2\theta\sin\theta - \sin^3\theta = 3\sin\theta - 4\sin^3\theta.}
$$

(On a utilisé $\sin^2\theta = 1 - \cos^2\theta$ et $\cos^2\theta = 1 - \sin^2\theta$.)


% ------------------------------------------------------------
%  NOTION 5 bis : Linéarisation et intégration
% ------------------------------------------------------------

\exerciceresolu
\textbf{Linéariser $\cos^4\theta$, puis calculer $\displaystyle\int_0^{\pi/2}\cos^4\theta\dd\theta$}

\solution

\textbf{1. Linéarisation} (méthode du résumé p.~\pageref{res:cplx_linear}).

\emph{Étape 1 — Euler.} $\cos\theta = \dfrac{\e^{\ii\theta}+\e^{-\ii\theta}}{2}$, donc
$$
\cos^4\theta = \frac{1}{2^4}\left(\e^{\ii\theta}+\e^{-\ii\theta}\right)^{4}.
$$

\emph{Étape 2 — binôme de Newton.} Avec $a = \e^{\ii\theta}$ et $b = \e^{-\ii\theta}$,
$ab = 1$ :
$$
(a+b)^4 = a^4 + 4a^3b + 6a^2b^2 + 4ab^3 + b^4
= \e^{4\ii\theta} + 4\e^{2\ii\theta} + 6 + 4\e^{-2\ii\theta} + \e^{-4\ii\theta}.
$$

\emph{Étape 3 — regrouper les conjugués.} $\e^{\ii k\theta} + \e^{-\ii k\theta} = 2\cos(k\theta)$ :
$$
(a+b)^4 = 2\cos(4\theta) + 8\cos(2\theta) + 6.
$$

D'où, en divisant par $16$ :
$$
\boxed{\cos^4\theta = \frac{\cos(4\theta) + 4\cos(2\theta) + 3}{8}.}
$$

\emph{Contrôle rapide :} pour $\theta = 0$, le membre de droite vaut
$\dfrac{1+4+3}{8} = 1 = \cos^4 0$. \checkmark

\textbf{2. Intégration.} La forme linéarisée s'intègre terme à terme :
$$
\int_0^{\pi/2}\cos^4\theta\dd\theta
= \frac{1}{8}\left[\frac{\sin(4\theta)}{4} + 4\cdot\frac{\sin(2\theta)}{2}
+ 3\theta\right]_0^{\pi/2}.
$$
Or $\sin(2\pi) = 0$ et $\sin(\pi) = 0$, donc il ne reste que le terme en $\theta$ :
$$
\int_0^{\pi/2}\cos^4\theta\dd\theta = \frac{1}{8}\times 3\times\frac{\pi}{2}
= \frac{3\pi}{16}.
$$

\begin{astucebac}
	C'est \emph{la} raison d'être de la linéarisation : $\cos^4\theta$ ne
	s'intègre pas tel quel, alors que $\cos(4\theta)$, $\cos(2\theta)$ et $3$
	s'intègrent immédiatement. Réflexe : \emph{une puissance de $\cos$ ou de
	$\sin$ sous une intégrale $\Rightarrow$ linéariser d'abord}.
\end{astucebac}


% ------------------------------------------------------------
%  NOTION 6 : Équation du second degré dans C
% ------------------------------------------------------------

\exerciceresolu
\textbf{Résoudre une équation du second degré dans $\C$}

Résoudre dans $\C$ l'équation $z^2 - (3 + \ii)z + (2 + 3\ii) = 0$.

\solution

On calcule le discriminant :
$$
\Delta = (3+\ii)^2 - 4(2+3\ii) = 9 + 6\ii - 1 - 8 - 12\ii = -6\ii.
$$

On cherche $\delta = x + y\ii$ tel que $\delta^2 = -6\ii$ :
$$
\begin{cases}
	x^2 - y^2 = 0 \\
	2xy = -6 \\
	x^2 + y^2 = \abs{-6\ii} = 6
\end{cases}
\implies
\begin{cases}
	x^2 = y^2 \\
	xy = -3 \\
	x^2 + y^2 = 6
\end{cases}
\implies x^2 = 3 \implies x = \sqrt{3}.
$$

Comme $xy = -3 < 0$, on a $y = -\sqrt{3}$, donc $\delta = \sqrt{3} - \sqrt{3}\,\ii$.

Les solutions sont :
$$
\begin{aligned}
	z_1 &= \frac{(3+\ii) + (\sqrt{3} - \sqrt{3}\,\ii)}{2}
	     = \frac{(3+\sqrt{3}) + (1-\sqrt{3})\ii}{2},\\
	z_2 &= \frac{(3+\ii) - (\sqrt{3} - \sqrt{3}\,\ii)}{2}
	     = \frac{(3-\sqrt{3}) + (1+\sqrt{3})\ii}{2}.
\end{aligned}
$$



%\exerciceresolu
%\textbf{Système somme–produit}
%
%Déterminer les complexes $z$ et $z'$ tels que $z + z' = 2 - 2\ii$ et
%$z\,z' = -2\ii$.
%
%\solution
%
%$z$ et $z'$ sont les solutions de $X^2 - (2 - 2\ii)X - 2\ii = 0$.
%
%$\Delta = (2 - 2\ii)^2 + 8\ii = (4 - 8\ii - 4) + 8\ii = 0$.
%
%Racine double : $X = \dfrac{2 - 2\ii}{2} = 1 - \ii$. Donc $z = z' = 1 - \ii$.

\exerciceresolu
\textbf{Ensemble de points}

Le plan complexe est rapporté à un repère orthonormé direct. On note $A$ le
point d'affixe $\ii$ et $B$ celui d'affixe $2$. Déterminer et construire
l'ensemble $(\Gamma)$ des points $M$ d'affixe $z$ ($z \neq \ii$) tels que
$$
Z = \frac{z - 2}{z - \ii} \quad \text{soit un imaginaire pur.}
$$

\solution

\textbf{Interprétation.} $\arg(Z) =\bigl(\vect{AM},\,\vect{BM}\bigr)\ [2\pi]$ n'est défini que si
$M \neq A, B$. $Z \in \ii\R^*$ signifie
$\bigl(\vect{MA},\,\vect{MB}\bigr) \equiv \dfrac{\pi}{2}\ [\pi]$, c'est-à-dire
$\vect{MA} \perp \vect{MB}$.

$(\Gamma)$ est donc le \textbf{cercle de diamètre $[AB]$}, privé des points $A$
et $B$. (Le cas $Z = 0$, soit $M = B$, est exclu ; $M = A$ est exclu car
$z \neq \ii$.)

\textbf{Contrôle analytique.} Avec $z = x + \ii y$ :
$$
Z = \frac{(x - 2) + \ii y}{x + \ii(y - 1)}
= \frac{\bigl[(x-2) + \ii y\bigr]\bigl[x - \ii(y-1)\bigr]}{x^2 + (y-1)^2}.
$$
$Z$ est imaginaire pur $\iff \mathrm{Re}(Z) = 0 \iff x(x - 2) + y(y - 1) = 0$,
soit $\left(x - 1\right)^2 + \left(y - \dfrac12\right)^2 = \dfrac54$ : cercle de
centre $\Omega\!\left(1\,;\dfrac12\right)$ (milieu de $[AB]$) et de rayon
$\dfrac{\sqrt5}{2} = \dfrac{AB}{2}$.


% ------------------------------------------------------------
%  NOTION 7 : Racines n-ièmes de l'unité
% ------------------------------------------------------------

\exerciceresolu
\textbf{Déterminer les racines $4$-ièmes de l'unité et les placer dans le plan}

Résoudre $z^4 = 1$ dans $\C$.

\solution

Les racines $4$-ièmes de l'unité sont $z_k = \e^{\ii\frac{2k\pi}{4}} = \e^{\ii\frac{k\pi}{2}}$
pour $k \in \{0, 1, 2, 3\}$ :

$$
z_0 = \e^{0} = 1, \quad
z_1 = \e^{\ii\frac{\pi}{2}} = \ii, \quad
z_2 = \e^{\ii\pi} = -1, \quad
z_3 = \e^{\ii\frac{3\pi}{2}} = -\ii.
$$

\textbf{Vérification :} $1^4 = 1$ \checkmark, $\ii^4 = (\ii^2)^2 = (-1)^2 = 1$ \checkmark,
$(-1)^4 = 1$ \checkmark, $(-\ii)^4 = 1$ \checkmark.

Ces quatre points sont les sommets du carré inscrit dans le cercle unité.

\remarque
La somme des racines $n$-ièmes de l'unité est toujours nulle :
$1 + \ii + (-1) + (-\ii) = 0$. \checkmark


% ------------------------------------------------------------
%  NOTION 8 : Géométrie — affixe, distance, milieu
% ------------------------------------------------------------

\exerciceresolu
\textbf{Utiliser les complexes en géométrie}

On donne $A$ d'affixe $z_A = 2 + 3\ii$, $B$ d'affixe $z_B = -1 + \ii$
et $C$ d'affixe $z_C = 3 - \ii$.

\begin{enumerate}
	\item Calculer $AB$, $BC$ et $AC$.
	\item Déterminer l'affixe du milieu $I$ de $[AC]$.
	\item Montrer que le triangle $ABC$ est rectangle en $B$.
\end{enumerate}

\solution

\textbf{1.}
$$
\begin{aligned}
	AB &= \abs{z_B - z_A} = \abs{(-1+\ii)-(2+3\ii)} = \abs{-3-2\ii} = \sqrt{13},\\
	BC &= \abs{z_C - z_B} = \abs{(3-\ii)-(-1+\ii)} = \abs{4-2\ii} = 2\sqrt{5},\\
	AC &= \abs{z_C - z_A} = \abs{(3-\ii)-(2+3\ii)} = \abs{1-4\ii} = \sqrt{17}.
\end{aligned}
$$

\textbf{2.} L'affixe du milieu $I$ est :
$$
z_I = \frac{z_A + z_C}{2} = \frac{(2+3\ii)+(3-\ii)}{2} = \frac{5+2\ii}{2} = \frac{5}{2} + \ii.
$$

\textbf{3.} On calcule $\dfrac{z_C - z_B}{z_A - z_B}$ :
$$
\frac{z_C - z_B}{z_A - z_B}
= \frac{4 - 2\ii}{3 + 2\ii}
= \frac{(4-2\ii)(3-2\ii)}{(3+2\ii)(3-2\ii)}
= \frac{12 - 8\ii - 6\ii + 4\ii^2}{9+4}
= \frac{8 - 14\ii}{13}
= \frac{8}{13} - \frac{14}{13}\ii.
$$

Ce rapport n'est pas imaginaire pur, donc $(BC)$ n'est pas perpendiculaire à $(BA)$.

Vérifions par le théorème de Pythagore :
$$
\begin{aligned}
	AB^2 + BC^2 &= 33 \neq 17 = AC^2,\\
	AB^2 + AC^2 &= 30 \neq 20 = BC^2,\\
	BC^2 + AC^2 &= 37 \neq 13 = AB^2.
\end{aligned}
$$

Le triangle $ABC$ \textbf{n'est pas rectangle}.

\begin{attention}
	Pour montrer qu'un triangle est rectangle en $B$, on vérifie que
	$\dfrac{z_C - z_B}{z_A - z_B}$ est \textbf{imaginaire pur}
	(partie réelle nulle), ou que $AB^2 + BC^2 = AC^2$.
\end{attention}


% ------------------------------------------------------------
%  NOTION 9 : Argument — alignement et perpendicularité
% ------------------------------------------------------------

\exerciceresolu
\textbf{Utiliser les arguments pour étudier une configuration géométrique}

Soient $A$, $B$, $C$, $D$ d'affixes respectives $\ii$, $2+\ii$, $1+3\ii$, $-1+2\ii$.

\begin{enumerate}
	\item Montrer que $A$, $B$, $C$ ne sont pas alignés.
	\item Montrer que $(AB) \perp (AD)$.
\end{enumerate}

\solution

\textbf{1.} On calcule $\dfrac{z_C - z_A}{z_B - z_A}$ :
$$
\frac{z_C - z_A}{z_B - z_A}
= \frac{(1+3\ii) - \ii}{(2+\ii) - \ii}
= \frac{1 + 2\ii}{2}
= \frac{1}{2} + \ii.
$$
Ce quotient n'est pas réel (partie imaginaire $= 1 \neq 0$),
donc $A$, $B$, $C$ \textbf{ne sont pas alignés}.

\textbf{2.} On calcule $\dfrac{z_D - z_A}{z_B - z_A}$ :
$$
\frac{z_D - z_A}{z_B - z_A}
= \frac{(-1+2\ii) - \ii}{(2+\ii) - \ii}
= \frac{-1 + \ii}{2}
= -\frac{1}{2} + \frac{1}{2}\ii.
$$
Ce quotient est imaginaire pur si et seulement si sa partie réelle est nulle.
Ici, la partie réelle est $-\dfrac{1}{2} \neq 0$, donc $(AB)$ et $(AD)$ ne sont pas perpendiculaires.

\begin{attention}
	$(AB) \perp (CD) \iff \dfrac{z_D - z_C}{z_B - z_A}$ est \textbf{imaginaire pur},
	c'est-à-dire que sa partie réelle est nulle.
\end{attention}


% ------------------------------------------------------------
%  NOTION 10 : Similitudes directes
% ------------------------------------------------------------

\exerciceresolu
\textbf{Identifier et caractériser une similitude directe}

Soit $f$ la transformation du plan d'expression complexe $z' = (1-\ii)z + 2 + 3\ii$.

\begin{listecol}(2)
	\task Identifier la nature de $f$.
	\task Déterminer le rapport et l'angle de $f$.
	\task Déterminer le centre de $f$.
\end{listecol}

\solution

\textbf{1.} $f$ est de la forme $z' = az + b$ avec $a = 1 - \ii \neq 1$ et $b = 2 + 3\ii$.
C'est une \textbf{similitude directe} (qui n'est ni une translation, ni une rotation pure,
ni une homothétie pure).

\textbf{2.}
$$
\begin{aligned}
	\text{Rapport} : r      &= \abs{a} = \abs{1-\ii} = \sqrt{2},\\
	\text{Angle} : \theta   &= \arg(1-\ii) = -\frac{\pi}{4}.
\end{aligned}
$$
$f$ est une similitude directe de rapport $\sqrt{2}$ et d'angle $-\dfrac{\pi}{4}$.

\textbf{3.} Le centre $\Omega$ d'affixe $\omega$ est le point fixe de $f$ :
$$
\omega = a\omega + b \implies \omega(1-a) = b
\implies \omega = \frac{b}{1-a} = \frac{2+3\ii}{1-(1-\ii)} = \frac{2+3\ii}{\ii}.
$$

On simplifie en multipliant par $\dfrac{\bar{\ii}}{\bar{\ii}} = \dfrac{-\ii}{-\ii}$ :
$$
\omega = \frac{(2+3\ii)(-\ii)}{(-\ii)(\ii)} = \frac{-2\ii - 3\ii^2}{1} = 3 - 2\ii.
$$

Le centre de $f$ est le point $\Omega$ d'affixe $3 - 2\ii$.

\begin{astucebac}
	Pour identifier rapidement $f : z \mapsto az + b$ :
	\begin{itemize}[label=\textbullet]
		\item $a = 1$ : translation de vecteur d'affixe $b$.
		\item $\abs{a} = 1$, $a \neq 1$ : rotation de centre $\dfrac{b}{1-a}$ et d'angle $\arg(a)$.
		\item $a \in \R^*$, $a \neq 1$ : homothétie de centre $\dfrac{b}{1-a}$ et de rapport $a$.
		\item Sinon : similitude directe de rapport $\abs{a}$, d'angle $\arg(a)$
		et de centre $\dfrac{b}{1-a}$.
	\end{itemize}
\end{astucebac}



\section{Exercices d'entraînement}

\subsection{Échauffement : lire une figure}

\exercice{Placer, lire, calculer \hfill \facile}
Dans le plan complexe, on donne $A$ d'affixe $2 + \ii$, $B$ d'affixe
$-1 + 3\ii$ et $C$ d'affixe $3 - 2\ii$.
\begin{enumerate}
	\item Placer $A$, $B$, $C$.
	\item Calculer l'affixe de $\vect{AB}$ et la distance $AB$.
	\item Déterminer l'affixe du milieu $I$ de $[BC]$ et celle du point $D$ tel
	que $ABCD$ soit un parallélogramme.
\end{enumerate}

\exercice{Module et argument sur la figure \hfill \facile}
On considère les points d'affixes $z_1 = 1 + \ii$, $z_2 = -2$, $z_3 = -\ii$,
$z_4 = \sqrt3 + \ii$.
\begin{enumerate}
	\item Sans calcul écrit, placer ces quatre points et lire graphiquement
	leur module et un argument.
	\item Écrire chacun sous forme exponentielle.
	\item Placer le point d'affixe $z_1 z_4$ ; que constate-t-on pour son
	argument ?
\end{enumerate}

\exercice{Reconnaître un ensemble simple \hfill \facile}
Associer chaque condition à sa représentation (cercle, droite, médiatrice,
demi-droite, cercle de diamètre) :
\begin{listecol}(2)
	\task $\abs{z - 1 + \ii} = 3$
	\task $\abs{z - 2} = \abs{z + \ii}$
	\task $\mathrm{Re}(z) = -1$
	\task $\arg(z - \ii) \equiv \pi/4\ [2\pi]$
	\task $\frac{z - 1}{z + 1} \in \ii\R$
\end{listecol}

\exercice{Nature d'un triangle \hfill \moyen}
$A$, $B$, $C$ ont pour affixes $z_A = 1$, $z_B = 1 + \ii$, $z_C = 2 + \ii$.
\begin{enumerate}
	\item Calculer $Z = \dfrac{z_C - z_A}{z_B - z_A}$.
	\item En déduire la nature du triangle $ABC$.
	\item Retrouver le résultat en calculant les longueurs $AB$, $AC$, $BC$.
\end{enumerate}

\exercice{Vrai ou faux \hfill \facile}
Répondre par \emph{vrai} ou \emph{faux} en justifiant brièvement.
\begin{enumerate}
	\item Si $\abs{z} = \abs{z'}$ alors $z = z'$ ou $z = -z'$.
	\item $\arg(z\,z') \equiv \arg(z) + \arg(z')\ [2\pi]$ pour tous $z, z'$ non nuls.
	\item Une équation de degré $3$ à coefficients réels a toujours au moins une
	racine réelle.
	\item L'ensemble des points $M$ tels que $\abs{z - a} = \abs{z - b}$ est un
	cercle.
	\item $\dfrac{z - a}{z - b}$ est réel si et seulement si $M$ appartient à la
	droite $(AB)$.
	\item Si $z$ et $\bar z$ sont solutions de $z^2 + bz + c = 0$ alors $b$ et
	$c$ sont réels.
	\item Les images des racines $n$-ièmes de l'unité forment un polygone
	régulier.
	\item $\bigl(\e^{\ii\theta}\bigr)^{n} = \e^{\ii n\theta}$ seulement si
	$n$ est un entier positif.
	\item $\e^{\ii\theta} = \e^{\ii\theta'}$ équivaut à $\theta = \theta'$.
	\item Pour tout réel $\theta$, $\cos^2\theta = \dfrac{1+\cos(2\theta)}{2}$.
\end{enumerate}

\subsection{Formes, équations et systèmes}

\exercice{Calculs et formes \hfill \facile}
\begin{enumerate}
	\item Écrire $\dfrac{(1 + \ii)^3}{2 - \ii}$ sous forme algébrique.
	\item Écrire $z = -\sqrt3 + \ii$ sous forme exponentielle, puis calculer
	$z^{6}$.
	\item Résoudre dans $\C$ : $\bar z = 2\ii z - 1$.
\end{enumerate}

\exercice{Second degré et racine carrée \hfill \moyen}
\begin{enumerate}
	\item Déterminer les racines carrées de $-3 + 4\ii$.
	\item Résoudre $z^2 - (3 + 2\ii)z + (1 + 3\ii) = 0$.
	\item Résoudre $z^2 - 2z\cos\theta + 1 = 0$ ($\theta \in \R$) et interpréter
	les solutions sur le cercle unité.
\end{enumerate}

\exercice{Degré 3 et 4 \hfill \moyen}
\begin{enumerate}
	\item Résoudre $z^3 - 1 = 0$, puis $z^3 = 8\ii$.
	\item Résoudre $z^4 + z^2 + 1 = 0$ (poser $Z = z^2$).
	\item Soit $P(z) = z^4 - 2z^3 + 2z^2 + 2z - 3$. Vérifier que $\ii$ et
	$-\ii$ sont racines, puis résoudre $P(z) = 0$.
\end{enumerate}

\exercice{Systèmes \hfill \moyen}
\begin{enumerate}
	\item Résoudre $\begin{cases} z + z' = 1 + 3\ii \\ z\,z' = -2 + 2\ii \end{cases}$.
	\item Résoudre $\begin{cases} 2z + \ii z' = 3 \\ z - z' = 1 - \ii \end{cases}$.
	\item Déterminer $z$ tel que $z + 2\bar z = 3 - 6\ii$.
\end{enumerate}

\exercice{Forme exponentielle : calculer vite \hfill \facile}
\emph{Résumé p.~\pageref{res:cplx_expo}.}
\begin{enumerate}
	\item Écrire sous forme exponentielle : $-4$, \ $3\ii$, \ $1 - \ii$, \
	$-\sqrt3 + 3\ii$.
	\item Soit $z = 2\e^{\ii\pi/3}$ et $z' = \sqrt2\,\e^{-\ii\pi/4}$.
	Donner la forme exponentielle de $zz'$, $\dfrac{z}{z'}$, $\bar z$,
	$\dfrac{1}{z}$ et $-z$.
	\item En déduire la forme algébrique de $zz'$.
\end{enumerate}

\exercice{Puissances et formule de Moivre \hfill \moyen}
\emph{Résumé p.~\pageref{res:cplx_moivre}.}
\begin{enumerate}
	\item Calculer $(1 - \ii)^{10}$ et $\left(-1 + \ii\sqrt3\right)^{6}$.
	\item Exprimer $\sin(4x)$ en fonction de $\cos x$ et $\sin x$.
	\item Soit $z = \e^{\ii\pi/5}$. Montrer que $z^{10} = 1$, puis déterminer
	le plus petit entier $n \geq 1$ tel que $z^n$ soit réel.
\end{enumerate}

\exercice{Trigonométrie \hfill \difficile}
\emph{Résumés p.~\pageref{res:cplx_moivre} et p.~\pageref{res:cplx_linear}.}
\begin{enumerate}
	\item Linéariser $\cos^3 x$ et $\sin^4 x$.
	\item En déduire $\displaystyle\int_0^{\pi/2}\sin^4 x\dd x$.
	\item Exprimer $\cos(3x)$ en fonction de $\cos x$.
	\item Écrire $\cos x + \sqrt3\,\sin x$ sous la forme $R\cos(x - \varphi)$,
	puis résoudre $\cos x + \sqrt3\,\sin x = 1$.
\end{enumerate}

\subsection{Géométrie et ensembles de points}

\exercice{Alignement et orthogonalité \hfill \moyen}
$A$, $B$, $C$, $D$ ont pour affixes $1 + 2\ii$, $3 - \ii$, $-\ii$, $4 + 4\ii$.
\begin{enumerate}
	\item Les points $A$, $B$, $C$ sont-ils alignés ?
	\item Les droites $(AB)$ et $(CD)$ sont-elles perpendiculaires ?
	\item Calculer une mesure de l'angle $\bigl(\vect{CA},\,\vect{CB}\bigr)$.
\end{enumerate}

\exercice{Ensembles définis par un module \hfill \moyen}
On note $A(1 - \ii)$ et $B(3 + \ii)$. Déterminer et construire l'ensemble des
points $M$ d'affixe $z$ tels que :
\begin{listecol}(2)
	\task $\abs{z - 1 + \ii} = 2$
	\task $\abs{z - 1 + \ii} = \abs{z - 3 - \ii}$
	\task $\abs{z - 1 + \ii} = 2\,\abs{z - 3 - \ii}$
	\task $\abs{2z - 4} = \abs{z - 3 - \ii}$
\end{listecol}

\exercice{Ensembles définis par un argument ou un quotient \hfill \difficile}
On note $A$ le point d'affixe $1$ et $B$ celui d'affixe $\ii$. Pour
$z \notin \{1, \ii\}$, on pose $Z = \dfrac{z - 1}{z - \ii}$.
\begin{enumerate}
	\item Déterminer l'ensemble $(E_1)$ des points $M$ tels que $Z \in \R$.
	\item Déterminer l'ensemble $(E_2)$ des points $M$ tels que $Z \in \ii\R$.
	\item Déterminer l'ensemble $(E_3)$ des points $M$ tels que $\abs{Z} = 1$.
	\item Déterminer l'ensemble $(E_4)$ des points $M$ tels que
	$\arg(Z) \equiv \dfrac{\pi}{2}\ [2\pi]$ et le comparer à $(E_2)$.
\end{enumerate}

\exercice{Un ensemble par la partie réelle \hfill \difficile}
Pour $z \neq -\ii$, on pose $Z = \dfrac{z - 1}{z + \ii}$.
\begin{enumerate}
	\item Écrire $Z$ sous forme algébrique en posant $z = x + \ii y$.
	\item Déterminer l'ensemble des points $M$ tels que $Z$ soit réel.
	\item Déterminer l'ensemble des points $M$ tels que $\abs{Z} = 1$.
	\item Déterminer l'ensemble des points $M$ tels que $\mathrm{Re}(Z) = 1$.
\end{enumerate}

\subsection{Transformations}
% ============================================================
%   EXERCICES — Nature et éléments caractéristiques
%                d'une application du plan complexe
%   Mazava Maths
% ============================================================


% ------------------------------------------------------------
\exercice{Identifier une translation ou une rotation \hfill \etoilepleine\etoilepleine\etoilevide}
% ------------------------------------------------------------

Pour chacune des applications suivantes, déterminer la nature et les éléments
caractéristiques.

\begin{listecol}(2)
	\task $f_1 : z \mapsto z + 3 - 2\ii$
	\task $f_2 : z \mapsto \ii z + (1 - \ii)$
	\task $f_3 : z \mapsto -z + 4 + 2\ii$
	\task $f_4 : z \mapsto \e^{\ii\pi/3} z + 2$
\end{listecol}


% ------------------------------------------------------------
\exercice{Similitude directe — cas général \hfill \etoilepleine\etoilepleine\etoilevide}
% ------------------------------------------------------------

Soit $f$ l'application d'expression complexe
$$
z' = (1 + \ii\sqrt{3})\,z - 2 + 2\ii\sqrt{3}.
$$

\begin{enumerate}
	\item Montrer que $f$ est une similitude directe.
	\item Déterminer le rapport $k$ et l'angle $\theta$ de $f$.
	\item Déterminer l'affixe du centre $\Omega$ de $f$.
	\item Vérifier que $\Omega$ est bien un point invariant  de $f$.
\end{enumerate}


% ------------------------------------------------------------
\exercice{Déterminer une application à partir de deux images \hfill \etoilepleine\etoilepleine\etoilevide}
% ------------------------------------------------------------

On cherche une similitude directe $f : z' = az + b$ telle que
$$
f(0) = 2 + \ii \quad \text{et} \quad f(1) = 3 + 3\ii.
$$

\begin{enumerate}
	\item Déterminer $a$ et $b$.
	\item Préciser la nature de $f$, son rapport et son angle.
	\item Déterminer le centre de $f$ si ce n'est pas une translation.
\end{enumerate}


% ------------------------------------------------------------
\exercice{Composition de deux applications \hfill \etoilepleine\etoilepleine\etoilevide}
% ------------------------------------------------------------

On considère les applications
$$
f : z \mapsto \ii z + 1 + 2\ii \qquad \text{et} \qquad g : z \mapsto -\ii z + 3.
$$

\begin{enumerate}
	\item Préciser la nature de $f$ et de $g$.
	\item Déterminer l'expression complexe de $g \circ f$.
	\item Préciser la nature de $g \circ f$ et ses éléments caractéristiques.
	\item La composition $g \circ f$ est-elle commutative ? Justifier.
\end{enumerate}


% ------------------------------------------------------------
\exercice{Application définie par une condition géométrique \hfill \etoilepleine\etoilepleine\etoilevide}
% ------------------------------------------------------------

On considère la rotation $r$ de centre $\Omega$ d'affixe $1 + \ii$ et d'angle $\dfrac{\pi}{2}$.

\begin{enumerate}
	\item Écrire l'expression complexe de $r$.
	\item Calculer les images des points $A(2\,;\,0)$, $B(0\,;\,2)$ et $C(3\,;\,3)$.
	\item Montrer que $r$ conserve les distances, c'est-à-dire $\abs{r(z) - r(z')} = \abs{z - z'}$
	pour tous $z, z' \in \C$.
\end{enumerate}


% ------------------------------------------------------------
\exercice{Similitude indirecte \hfill \etoilepleine\etoilepleine\etoilepleine}
% ------------------------------------------------------------

On considère l'application $f : z \mapsto (2 + \ii)\bar{z} - 1 + 3\ii$.

\begin{enumerate}
	\item Montrer que $f$ n'est pas une similitude directe.
	\item Montrer que $f$ est une similitude \emph{indirecte} en vérifiant que $f$
	conserve les distances mais inverse l'orientation.
	\item Déterminer les points fixes de $f$ (en posant $\bar{z} = \overline{z}$
	et en séparant parties réelle et imaginaire).
\end{enumerate}


% ------------------------------------------------------------
\exercice{Retrouver une application à partir de ses éléments \hfill \etoilepleine\etoilepleine\etoilepleine}
% ------------------------------------------------------------

\begin{enumerate}
	\item Déterminer l'expression complexe de la rotation de centre $A$ d'affixe
	$2 - \ii$ et d'angle $\dfrac{3\pi}{4}$.

	\item Déterminer l'expression complexe de l'homothétie de centre $B$ d'affixe
	$-1 + 2\ii$ et de rapport $-3$.

	\item Déterminer l'expression complexe de la similitude directe de centre $C$
	d'affixe $\ii$, de rapport $\sqrt{2}$ et d'angle $\dfrac{\pi}{4}$.
\end{enumerate}


% ------------------------------------------------------------
\exercice{Problème de synthèse \hfill \etoilepleine\etoilepleine\etoilepleine}
% ------------------------------------------------------------

Dans le plan complexe, on considère le carré $ABCD$ direct de centre $\Omega$
tel que $A$ a pour affixe $z_A = 0$ et $B$ a pour affixe $z_B = 2$.

\begin{enumerate}
	\item Déterminer les affixes de $C$, $D$ et $\Omega$.

	\item Soit $r$ la rotation de centre $\Omega$ et d'angle $\dfrac{\pi}{2}$.
	Écrire l'expression complexe de $r$ et vérifier que $r(A) = B$.

	\item Soit $s$ la similitude directe transformant $A$ en $B$ et $B$ en $C$.
	\begin{enumerate}
		\item Déterminer l'expression complexe de $s$.
		\item Préciser la nature, le rapport, l'angle et le centre de $s$.
		\item Vérifier que le centre de $s$ coïncide avec $\Omega$.
	\end{enumerate}

	\item Calculer $s^2$, $s^3$ et $s^4$. Que remarque-t-on ?
\end{enumerate}

% ============================================================
%   EXERCICES — Applications du plan complexe
%   Nature, points images, points invariants
%   Mazava Maths
% ============================================================


% ------------------------------------------------------------
\exercice{QCM — Identifier la nature d'une application \hfill \etoilepleine\etoilepleine\etoilevide}
% ------------------------------------------------------------

\textit{Pour chaque question, une seule réponse est correcte.}

\begin{enumerate}
	\item L'application $f : z \mapsto \ii z + 1 - \ii$ est :
	\begin{qcm}(2)
		\task Translation de vecteur $1 - \ii$
		\task Rotation de centre $(1+\ii)/(1-\ii)$, angle $\pi/2$
		\task Rotation de centre $1$, angle $\pi/2$
		\task Homothétie de centre $\ii$, rapport $1$
	\end{qcm}

	\item L'application $g : z \mapsto 2z - 2 - 4\ii$ est :
	\begin{qcm}(2)
		\task Rotation de rapport $2$, angle $0$
		\task Homothétie de centre $2 + 4\ii$, rapport $2$
		\task Similitude de rapport $2$, angle $\pi$
		\task Translation de vecteur $-2 - 4\ii$
	\end{qcm}

	\item La similitude $h : z \mapsto (1+\ii)z + 2$ a pour rapport :
	\begin{qcm}(4)
		\task $1$
		\task $\sqrt{2}$
		\task $2$
		\task $1 + \ii$
	\end{qcm}

	\item Le point fixe de $f : z \mapsto (2+\ii)z - 1 + 3\ii$ a pour affixe :
	\begin{qcm}(2)
		\task $(-1+3\ii)/(1+\ii)$
		\task $-1 - 2\ii$
		\task $(1-3\ii)/(1-\ii)$
		\task $-1 + 2\ii$
	\end{qcm}
\end{enumerate}


% ------------------------------------------------------------
\exercice{Application laissant un point invariant \hfill \etoilepleine\etoilepleine\etoilevide}
% ------------------------------------------------------------

Soit $f$ une similitude directe telle que :
$$
f\text{ laisse invariant le point }A \text{ d'affixe } z_A = 1 + \ii
\quad \text{et} \quad
f(B) = C,
$$
où $B$ a pour affixe $z_B = 3 + \ii$ et $C$ a pour affixe $z_C = 1 + 3\ii$.

\begin{enumerate}
	\item Justifier que $f$ est de la forme $z' = a(z - z_A) + z_A$
	avec $a \in \C^*$.
	\item En utilisant l'image de $B$, déterminer $a$.
	\item En déduire la nature de $f$, son rapport et son angle.
	\item Calculer $f(D)$ où $D$ a pour affixe $z_D = 1 - \ii$.
\end{enumerate}


% ------------------------------------------------------------
\exercice{Déterminer $f$ par deux conditions \hfill \etoilepleine\etoilepleine\etoilevide}
% ------------------------------------------------------------

On cherche une similitude directe $f : z' = az + b$ vérifiant :
$$
f(A) = B \quad \text{et} \quad f(B) = C,
$$
où $z_A = 0$, $z_B = 2$ et $z_C = 2 + 2\ii$.

\begin{enumerate}
	\item Écrire le système de deux équations vérifié par $a$ et $b$.
	\item Résoudre ce système pour déterminer $a$ et $b$.
	\item Préciser la nature, le rapport, l'angle et le centre de $f$.
	\item Que vaut $f(C)$ ? En déduire le quadrilatère $ABCD$ où $D = f(C)$.
\end{enumerate}


% ------------------------------------------------------------
\exercice{QCM — Points invariants et images \hfill \etoilepleine\etoilepleine\etoilevide}
% ------------------------------------------------------------

\textit{Pour chaque question, une seule réponse est correcte.}

Soit $r$ la rotation de centre $\Omega$ d'affixe $1 + \ii$ et d'angle $\dfrac{\pi}{2}$.

\begin{enumerate}
	\item L'expression complexe de $r$ est :
	\begin{qcm}(2)
		\task $z' = \ii z + 1 + \ii$
		\task $z' = \ii z + 2$
		\task $z' = \ii z - 2$
		\task $z' = \ii z + (1-\ii)(1-\ii)$
	\end{qcm}

	\item L'image du point $A$ d'affixe $3 + \ii$ par $r$ est le point d'affixe :
	\begin{qcm}(4)
		\task $3 + \ii$
		\task $-1 + 3\ii$
		\task $1 + 3\ii$
		\task $3 - \ii$
	\end{qcm}

	\item Le seul point invariant par $r$ est :
	\begin{qcm}(2)
		\task L'origine $O$
		\task Le point d'affixe $2$
		\task Le point $\Omega$ d'affixe $1 + \ii$
		\task Aucun point invariant
	\end{qcm}

	\item $r^2$ est :
	\begin{qcm}(2)
		\task Rotation de centre $\Omega$, angle $\pi$
		\task Translation de vecteur $2\ii$
		\task L'identité
		\task Homothétie de rapport $-1$
	\end{qcm}
\end{enumerate}


% ------------------------------------------------------------
\exercice{Application définie par un point fixe et une image \hfill \etoilepleine\etoilepleine\etoilepleine}
% ------------------------------------------------------------

Soit $f$ une similitude directe qui laisse invariant le point $\Omega$ d'affixe
$\omega = 2 + \ii$ et qui transforme $A$ d'affixe $z_A = 4 + \ii$ en $B$ d'affixe
$z_B = 2 + 3\ii$.

\begin{enumerate}
	\item Montrer que $f$ s'écrit $z' - \omega = a(z - \omega)$ et déterminer $a$.
	\item Calculer $\abs{a}$ et $\arg(a)$. En déduire la nature exacte de $f$.
	\item Déterminer l'image du point $C$ d'affixe $z_C = 2 - \ii$.
	\item Montrer que $\Omega$, $A$ et $C$ sont alignés.
	En déduire une propriété géométrique de $f$.
\end{enumerate}


% ------------------------------------------------------------
\exercice{Trouver $f$ sachant $f^2 = \mathrm{id}$ \hfill \etoilepleine\etoilepleine\etoilepleine}
% ------------------------------------------------------------

Soit $f : z' = az + b$ une similitude directe telle que $f \circ f = \mathrm{id}$
(l'application identité) et $f \neq \mathrm{id}$.

\begin{enumerate}
	\item Montrer que $f \circ f = \mathrm{id}$ implique $a^2 = 1$ et $b(a+1) = 0$.
	\item En déduire les valeurs possibles de $a$.
	\item Que se passe-t-il si $a = 1$ ? Si $a = -1$ ?
	\item On suppose $a = -1$. Déterminer $f$ sachant que $f$ transforme
	le point $A$ d'affixe $2 + \ii$ en le point $B$ d'affixe $-1 + 3\ii$.
	\item Quelle est la nature de $f$ ? Déterminer son centre.
\end{enumerate}


% ------------------------------------------------------------
\exercice{Problème de synthèse — triangle et similitude \hfill \etoilepleine\etoilepleine\etoilepleine}
% ------------------------------------------------------------

Dans le plan complexe, on donne le triangle équilatéral direct $ABC$ avec
$z_A = 0$ et $z_B = 2$.

\begin{enumerate}
	\item Rappeler que la rotation $r$ de centre $A$ et d'angle $\dfrac{\pi}{3}$
	transforme $B$ en $C$. En déduire $z_C$.

	\item Soit $f$ la similitude directe telle que $f(A) = B$ et $f(B) = C$.
	\begin{enumerate}
		\item Déterminer l'expression complexe de $f$.
		\item Montrer que le centre $\Omega$ de $f$ est le centre du triangle $ABC$.
		\item Calculer $f(C)$ et interpréter géométriquement.
	\end{enumerate}

	\item Soit $g = f \circ f \circ f$. Calculer l'expression complexe de $g$.
	Quelle est la nature de $g$ ?

	\item En déduire que $f^3 = \mathrm{id}$. Que représente géométriquement
	cette propriété ?
\end{enumerate}


% ------------------------------------------------------------
\exercice{QCM de synthèse \hfill \moyen}
% ------------------------------------------------------------

\textit{Plusieurs réponses peuvent être correctes.}

\begin{enumerate}
	\item Une similitude directe $f : z' = az + b$ est une rotation si et seulement si :
	\begin{qcm}(4)
		\task $a \in \R$
		\task $\abs{a} = 1$
		\task $\arg(a) = 0$
		\task $\abs{a} = 1$, $a \neq 1$
	\end{qcm}

	\item Si $f$ laisse invariants deux points distincts $A$ et $B$, alors :
	\begin{qcm}(2)
		\task Nécessairement l'identité
		\task Une translation
		\task Une rotation d'angle $2\pi$
		\task L'identité ou la symétrie d'axe $(AB)$
	\end{qcm}

	\item La composition de deux rotations de même centre $\Omega$ et d'angles
	$\alpha$ et $\beta$ est :
	\begin{qcm}(2)
		\task Rotation de centre $\Omega$, angle $\alpha + \beta$
		\task L'identité si $\alpha + \beta \equiv 0\,[2\pi]$
		\task Une translation si $\alpha + \beta \equiv 0\,[2\pi]$
		\task Similitude de rapport $2$
	\end{qcm}

	\item Si $f : z \mapsto az + b$ avec $a \neq 0$ et $a \neq 1$, alors le centre de $f$ :
	\begin{qcm}(2)
		\task A pour affixe $b/(1-a)$
		\task Est le seul point fixe de $f$
		\task A pour affixe $b/(a-1)$
		\task N'existe pas pour une translation
	\end{qcm}
\end{enumerate}



% ============================================================
%   PROBLÈME — Style BAC américain
%   Construction d'une spirale de triangles par similitude
%   Mazava Maths
% ============================================================

\exercice{La spirale de triangles \hfill \etoilepleine\etoilepleine\etoilepleine}

\textit{Les quatre parties sont liées. Lire l'énoncé en entier avant de commencer.}

\bigskip

Dans le plan complexe rapporté à un repère orthonormé direct $(O, \vec{\imath}, \vec{\jmath})$,
on considère le point $A_0$ d'affixe $z_0 = 4$ et on définit la similitude directe
$$
f : z' = \frac{1+\ii}{2}\,z.
$$

On pose $A_n$ le point d'affixe $z_n = f^n(z_0)$ pour tout $n \in \N$,
où $f^n$ désigne la composée $n$ fois de $f$ avec elle-même.

\bigskip
\noindent\rule{\linewidth}{0.3pt}
\bigskip

\noindent\textbf{Partie A — Étude de la similitude $f$}

\begin{enumerate}
	\item Montrer que $f$ est une similitude directe de centre $O$,
	de rapport $r = \dfrac{\sqrt{2}}{2}$ et d'angle $\theta = \dfrac{\pi}{4}$.

	\item Justifier que $O$ est le seul point invariant par $f$.

	\item Calculer $f^2$, $f^4$ et $f^8$.
	Quelle est la nature de chacune de ces applications ?

	\item Montrer que $f^8$ est une homothétie. Préciser son centre et son rapport.
\end{enumerate}

\bigskip
\noindent\rule{\linewidth}{0.3pt}
\bigskip

\noindent\textbf{Partie B — La suite de points $(A_n)$}

\begin{enumerate}
	\item Calculer $z_1$, $z_2$, $z_3$ et $z_4$ sous forme algébrique.
	Placer les points $O$, $A_0$, $A_1$, $A_2$, $A_3$, $A_4$ dans le plan.

	\item Montrer que $z_n = 4 \times \left(\dfrac{1+\ii}{2}\right)^n$
	et mettre $z_n$ sous forme exponentielle $r_n\,\e^{\ii\theta_n}$.

	\item \begin{enumerate}
		\item Montrer que $\abs{z_n} = 4 \times \left(\dfrac{\sqrt{2}}{2}\right)^n$.
		\item Justifier que les points $A_n$ se rapprochent de $O$ quand $n \to +\infty$.
		\item Calculer l'angle $\left(\vect{OA_0}, \vect{OA_n}\right)$ et interpréter
		géométriquement la suite $(A_n)$.
	\end{enumerate}

	\item Les triangles $OA_nA_{n+1}$ sont-ils tous semblables ?
	Calculer le rapport $\dfrac{OA_{n+1}}{OA_n}$ et l'angle $\widehat{A_nOA_{n+1}}$.
\end{enumerate}

\bigskip
\noindent\rule{\linewidth}{0.3pt}
\bigskip

\noindent\textbf{Partie C — Construction de la spirale}

\begin{enumerate}
	\item On relie successivement $A_0 \to A_1 \to A_2 \to \cdots \to A_n \to O$.
	Calculer la longueur totale
	$$
	L_n = \sum_{k=0}^{n} A_k A_{k+1}
	$$
	en fonction de $n$, où $A_{n+1} = O$ pour le dernier segment.

	\textit{Indication : } $A_kA_{k+1} = \abs{z_{k+1} - z_k}$.

	\item Montrer que la série $\displaystyle\sum_{k=0}^{+\infty} A_k A_{k+1}$ converge
	et calculer sa somme $L$.

	\item \begin{enumerate}
		\item Montrer que les triangles $OA_nA_{n+1}$ sont rectangles isocèles.
		\item En déduire que tous les points $A_n$ appartiennent à des cercles
		de centre $O$ dont les rayons forment une suite géométrique.
		\item Construire soigneusement la spirale pour $n = 0, 1, 2, 3, 4, 5$
		avec une règle et un compas (ou en calculant les coordonnées).
	\end{enumerate}
\end{enumerate}

\bigskip
\noindent\rule{\linewidth}{0.3pt}
\bigskip

\noindent\textbf{Partie D — Généralisation}

On remplace $f$ par la similitude directe $g : z' = \rho\,\e^{\ii\alpha}\,z$
avec $0 < \rho < 1$ et $\alpha \in \left]0\,;\,2\pi\right[$.

\begin{enumerate}
	\item Montrer que la suite $(B_n)$ définie par $w_n = g^n(w_0)$
	avec $w_0 \in \C^*$ vérifie $w_n = w_0\,\rho^n\,\e^{\ii n\alpha}$.

	\item Montrer que les points $B_n$ s'enroulent autour de $O$ en formant
	une \textbf{spirale logarithmique} en justifiant que :
	$$
	\abs{w_n} = \abs{w_0}\,\rho^n \to 0 \quad \text{et} \quad
	\arg(w_n) = \arg(w_0) + n\alpha \to +\infty.
	$$

	\item Pour quelles valeurs de $\rho$ et $\alpha$ la spirale est-elle :
	\begin{enumerate}
		\item un cercle ?
		\item une demi-droite issue de $O$ ?
		\item une spirale qui se resserre en tournant dans le sens direct ?
	\end{enumerate}

	\item \textbf{Question ouverte.}
	Si $\rho > 1$, décrire qualitativement le comportement de la suite $(B_n)$.
	La longueur totale $\displaystyle\sum_{k=0}^{+\infty} B_k B_{k+1}$ converge-t-elle encore ?
\end{enumerate}


\subsection{Sujets types Bac}

\exercice{Sujet type Bac — forme exponentielle, Moivre et linéarisation \hfill \difficile}
\emph{Le plan complexe est rapporté à un repère orthonormé direct
$(O, \vec{u}, \vec{v})$.}

\textbf{Partie A} — \textit{forme exponentielle et puissances}

On pose $z_0 = 1 + \ii\sqrt3$.
\begin{enumerate}
	\item Calculer $\abs{z_0}$ et un argument de $z_0$, puis écrire $z_0$ sous
	forme exponentielle.
	\item En déduire la forme algébrique de $z_0^{\,6}$.
	\item Déterminer le plus petit entier $n \geq 1$ tel que $z_0^{\,n}$ soit
	un réel strictement positif.
	\item On note $A_k$ le point d'affixe $z_0^{\,k}$ pour $k \in \{0;1;2\}$.
	Placer $A_0$, $A_1$, $A_2$ et justifier que
	$\arg\!\left(\dfrac{z_0^{\,2}}{z_0}\right) \equiv \dfrac{\pi}{3}\ [2\pi]$.
\end{enumerate}

\textbf{Partie B} — \textit{formule de Moivre}

\begin{enumerate}
	\item Énoncer la formule de Moivre.
	\item En développant $(\cos\theta + \ii\sin\theta)^4$ par le binôme de
	Newton, montrer que
	$\cos(4\theta) = 8\cos^4\theta - 8\cos^2\theta + 1$.
	\item En déduire les solutions dans $[0\,;\,\pi]$ de l'équation
	$8\cos^4\theta - 8\cos^2\theta + 1 = 0$.
\end{enumerate}

\textbf{Partie C} — \textit{linéarisation et intégrale}

\begin{enumerate}
	\item Rappeler les formules d'Euler donnant $\cos\theta$ et
	$\sin\theta$.
	\item Linéariser $\sin^3\theta$, c'est-à-dire l'écrire sous la forme
	$a\sin\theta + b\sin(3\theta)$ avec $a$, $b$ réels à préciser.
	\item En déduire la valeur exacte de
	$\displaystyle I = \int_0^{\pi}\sin^3\theta\dd\theta$.
	\item Vérifier la cohérence du signe de $I$ à l'aide du signe de
	$\sin\theta$ sur $[0\,;\,\pi]$.
\end{enumerate}

\exercice{Sujet type Bac — équation de degré 3 et configuration \hfill \difficile}
Dans l'ensemble $\C$, on considère le polynôme $P(z) = z^3 + z - 10$.

\textbf{Partie A}
\begin{enumerate}
	\item Calculer $P(2)$.
	\item Déterminer les réels $a$, $b$ tels que
	$P(z) = (z - 2)(z^2 + az + b)$.
	\item Résoudre dans $\C$ l'équation $P(z) = 0$.
\end{enumerate}

\textbf{Partie B} — \textit{le plan complexe est rapporté à un repère
orthonormé direct. $A$, $B$, $C$ sont les points d'affixes les solutions de
$P(z) = 0$, $z_A = 2$ étant réel.}
\begin{enumerate}
	\item Placer $A$, $B$, $C$.
	\item Calculer $\dfrac{z_C - z_A}{z_B - z_A}$ et en déduire la nature du
	triangle $ABC$.
	\item Déterminer et construire l'ensemble des points $M$ tels que
	$\abs{z - z_B} = \abs{z - z_C}$.
\end{enumerate}

\exercice{Sujet type Bac — forme exponentielle et ensemble de points \hfill \difficile}
On donne les points $A$ et $B$ d'affixes $z_A = \ii$ et $z_B = -\ii$. Pour tout
point $M$ d'affixe $z \neq \ii$, on pose
$$
Z = \frac{z + \ii}{z - \ii}.
$$

\textbf{Partie A}
\begin{enumerate}
	\item Interpréter géométriquement $\abs{Z}$ et $\arg(Z)$.
	\item Déterminer et construire l'ensemble $(\Gamma_1)$ des points $M$ tels
	que $\abs{Z} = 1$.
	\item Déterminer et construire l'ensemble $(\Gamma_2)$ des points $M$ tels
	que $Z$ soit un réel strictement négatif.
	\item Déterminer l'ensemble $(\Gamma_3)$ des points $M$ tels que $Z$ soit
	un imaginaire pur.
\end{enumerate}

\textbf{Partie B}
\begin{enumerate}
	\item On pose $z = x + \ii y$. Écrire $\mathrm{Re}(Z)$ et $\mathrm{Im}(Z)$
	en fonction de $x$ et $y$.
	\item Retrouver la nature de $(\Gamma_1)$ et $(\Gamma_3)$ par le calcul.
\end{enumerate}

\exercice{Sujet type Bac — module, argument et lieu \hfill \difficile}
Le plan complexe est rapporté à un repère orthonormé direct. On considère
l'application $f$ qui à tout point $M$ d'affixe $z \neq 1$ associe le point
$M'$ d'affixe $z' = \dfrac{z + 1}{z - 1}$.

\textbf{Partie A}
\begin{enumerate}
	\item Déterminer les points invariants de $f$.
	\item Montrer que $z'$ est réel si et seulement si $M$ appartient à une
	droite à préciser.
	\item Montrer que $z'$ est imaginaire pur si et seulement si $M$ appartient
	à un cercle privé de deux points, à préciser.
\end{enumerate}

\textbf{Partie B}
\begin{enumerate}
	\item Déterminer l'ensemble des points $M$ tels que $\abs{z'} = 1$.
	\item Déterminer l'ensemble des points $M$ tels que $\abs{z'} = 2$.
\end{enumerate}

\exercice{Sujet type Bac — équation de degré 4 \hfill \difficile}
On considère l'équation $(E) : z^4 - 4z^3 + 12z^2 - 16z + 16 = 0$.

\textbf{Partie A}
\begin{enumerate}
	\item Vérifier que $(E)$ s'écrit $(z^2 - 2z + 4)^2 = 0$.
	\item En déduire les solutions de $(E)$ et les mettre sous forme
	exponentielle.
\end{enumerate}

\textbf{Partie B} — \textit{$A$ et $B$ sont les points d'affixes
$z_A = 1 + \ii\sqrt3$ et $z_B = 1 - \ii\sqrt3$ ; $O$ est l'origine.}
\begin{enumerate}
	\item Calculer $OA$, $OB$, $AB$ et une mesure de l'angle
	$\bigl(\vect{OA},\vect{OB}\bigr)$. En déduire la nature de $OAB$.
	\item Déterminer l'affixe du point $C$ tel que $OACB$ soit un losange.
	Ce losange est-il un carré ? Justifier.
	\item Déterminer et construire l'ensemble des points $M$ d'affixe $z$ tels
	que $\abs{z - z_A} = \abs{z - z_B}$.
\end{enumerate}

\exercice{Sujet type Bac — système, rotation et ensemble \hfill \difficile}
Le plan complexe est rapporté à un repère orthonormé direct. On note $A$, $B$
les points d'affixes $z_A = 1 + \ii$ et $z_B = -1 + 3\ii$.

\textbf{Partie A}
\begin{enumerate}
	\item Déterminer les complexes $z$ et $z'$ vérifiant
	$$
	\begin{cases}
		z + z' = 2\ii \\
		z\,z' = -2
	\end{cases}
	$$
	\item Soit $r$ la rotation de centre $A$ et d'angle $\dfrac{\pi}{2}$.
	Déterminer l'écriture complexe de $r$, puis l'affixe de $B' = r(B)$.
	\item Montrer que le triangle $ABB'$ est rectangle isocèle en $A$.
\end{enumerate}

\textbf{Partie B}
\begin{enumerate}
	\item Déterminer et construire l'ensemble $(\mathscr{C})$ des points $M$
	d'affixe $z$ tels que $\abs{z - 1 - \ii} = \abs{z + 1 - 3\ii}$.
	\item Déterminer et construire l'ensemble $(\mathscr{D})$ des points $M$
	tels que $\dfrac{z - z_A}{z - z_B}$ soit un imaginaire pur.
	\item Préciser l'intersection de $(\mathscr{C})$ et $(\mathscr{D})$.
\end{enumerate}
